%====================================================================================
%  DOCUMENT PRINCIPAL - TEMPLATE RAPPORT & MÉMOIRE UNIVERSITAIRE HAUT DE GAMME
%  Guide Pratique et Démonstration Exhaustive de ClasseGuillaume
%  Moteur recommandé : LuaLaTeX
%  Compilation : latexmk -lualatex -outdir=build main.tex
%====================================================================================
\documentclass[
  % draft,             % Activer pour masquer les images et accélérer la compilation
  french,              % Langue principale
  english,             % Langue secondaire
  font=calibri,        % Options de police : calibri, heros, libertinus, comicsans, standard
  todoNote             % Activer l'affichage des notes et révisions
]{ClasseGuillaume}

% --- PAQUETS COMPLÉMENTAIRES DIDACTIQUES ---
\RequirePackage{multicol}      % Mises en page dynamiques en 2 et 3 colonnes
\RequirePackage{pdflscape}     % Rotation physique des pages en mode paysage
\RequirePackage{tikz}          % Schémas vectoriels natifs

% --- MÉTADONNÉES DU DOCUMENT ---
\title{Le Guide Ultime de la Rédaction Scientifique et Académique avec \LaTeX{}}
\author{Guillaume \textsc{Winant}}
\date{\today}

% Robustesse des acronymes de moteurs dans les citations
\providecommand{\XeLaTeX}{Xe\LaTeX}

% --- BASE BIBLIOGRAPHIQUE ---
\addbibresource{bib/references.bib}

%====================================================================================
%  DÉBUT DU DOCUMENT
%====================================================================================
\begin{document}

%------------------------------------------------------------------------------------
%  PRÉ-CONTENU (FRONTMATTER : Numérotation romaine, sans filigrane)
%------------------------------------------------------------------------------------
\frontmatter
\masquerLogo % Le logo ne s'affiche pas sur la couverture ni les préliminaires

% 1. Page de garde principale (en géométrie large)
\geometrieLarge
%------------------------------------------------------------------------------------
%  PAGE DE GARDE PRINCIPALE HAUT DE GAMME
%  Utilise \geometrieLarge et \masquerLogo pour une couverture aérée et équilibrée
%------------------------------------------------------------------------------------
\begin{titlepage}
  \begin{singlespace}
  \centering

  % En-tête institutionnel avec logos
  \begin{minipage}[c]{0.28\textwidth}
    \IfFileExists{images/UCLouvainLogo.png}{%
      \includegraphics[height=1.8cm,keepaspectratio]{images/UCLouvainLogo.png}%
    }{}%
  \end{minipage}
  \hfill
  \begin{minipage}[c]{0.40\textwidth}
    \centering
    {\scshape\Large \textbf{UCLouvain}}\\[2mm]
    {\scshape\footnotesize Secteur des Sciences de la Santé}\\[1mm]
    {\scshape\footnotesize Faculté de Médecine et Médecine Dentaire}
  \end{minipage}
  \hfill
  \begin{minipage}[c]{0.28\textwidth}
    \raggedleft
    \IfFileExists{images/FacMedLogo.jpg}{%
      \includegraphics[height=1.8cm,keepaspectratio]{images/FacMedLogo.jpg}%
    }{}%
  \end{minipage}

  \vspace{1.8cm}

  % Type de livrable
  {\large\sffamily\bfseries\color{secondaryColor} RAPPORT TECHNIQUE ET GUIDE DE RÉFÉRENCE ACADÉMIQUE\par}
  \vspace{0.4cm}
  {\normalsize\scshape Modèle Typographique \& Démonstration Exhaustive de \texttt{ClasseGuillaume}\par}

  \vspace{1.2cm}

  % Cartouche de titre principal
  \begin{tcolorbox}[
    colframe=primaryColor,
    colback=primaryColor!4,
    boxrule=1.5pt,
    arc=6pt,
    left=12pt,
    right=12pt,
    top=14pt,
    bottom=14pt,
    halign=center
  ]
    {\Huge\sffamily\bfseries\color{primaryColor} \thetitle \par}
    \vspace{0.5cm}
    {\large\sffamily\color{darkgray} Manuel pratique, vitrine typographique et guide méthodologique complet pour la rédaction scientifique sous Lua\LaTeX{}\par}
  \end{tcolorbox}

  \vspace{1.5cm}

  % Cartouche d'auteurs et superviseurs
  \begin{minipage}[t]{0.46\textwidth}
    \begin{flushleft}\large
      \textcolor{primaryColor}{\bfseries\sffamily Auteur :}\\
      \textbf{\theauthor}\\
      \small Ingénieur \& Chercheur\\
      \footnotesize\texttt{guillaume.winant@student.uclouvain.be}
    \end{flushleft}
  \end{minipage}
  \hfill
  \begin{minipage}[t]{0.46\textwidth}
    \begin{flushright}\large
      \textcolor{primaryColor}{\bfseries\sffamily Direction Académique :}\\
      \textbf{Prof. Dr. Promoteur Principal}\\
      \small Chaire d'Ingénierie \& Typographie\\
      \footnotesize\texttt{superviseur@uclouvain.be}
    \end{flushright}
  \end{minipage}

  \vfill

  % Bas de page institutionnel
  {\small\textit{Document rédigé dans le cadre des travaux de recherche et développement méthodologiques}\par}
  \vspace{0.2cm}
  {\footnotesize Département de Recherche en Sciences Appliquées et Médicales\par}

  \vspace{0.8cm}

  {\large\sffamily\bfseries Année Académique 2025--2026 \quad---\quad \thedate\par}

  \end{singlespace}
\end{titlepage}


% 2. Bascule en géométrie Word pour les préliminaires rédigés
\geometrieWord

% 3. Tableau de bord des révisions (avec cible cliquable pour \todoH)
\begin{singlespace}
\listoftodos[Tableau de Bord des Révisions et Tâches de Rédaction]
\label{table_todos}
\end{singlespace}

% 4. Résumé et Abstract bilingues avec mots-clés
%------------------------------------------------------------------------------------
%  RÉSUMÉ & ABSTRACT BILINGUE AVEC MOTS-CLÉS
%  Utilisation de \addchap pour une présence dans la table des matières sans numéro
%------------------------------------------------------------------------------------
\addchap{Résumé}

Ce document constitue à la fois un manuel de référence approfondi sur la typographie scientifique avec \LaTeX{} moderne et la vitrine exhaustive de la classe personnalisée \texttt{ClasseGuillaume}.
L'objectif central est de démontrer qu'il est possible d'allier une rigueur académique sans compromis à une esthétique visuelle équilibrée et aérée.
Toutes les facettes d'un mémoire ou d'un rapport d'ingénierie contemporain sont abordées : architecture modulaire des fichiers, compilation reproductible sous le moteur Lua\LaTeX{}, détection dynamique des polices de caractères, gestion fine des interlignes et élimination systématique des ruptures de lignes disgracieuses par microtypographie.

Le lecteur découvrira une mise en pratique détaillée des fonctionnalités exclusives de la classe :
le système collaboratif d'annotations marginales \texttt{\textbackslash todoH} avec troncature automatique et infobulles interactives, la gestion à la volée de géométries multiples (\texttt{\textbackslash geometrieWord}, \texttt{\textbackslash geometrieLarge}, \texttt{\textbackslash geometrieStandard}), la séparation stricte entre figures techniques et illustrations iconographiques dans des tables indépendantes, ainsi que l'harmonisation complète des tableaux sans barres verticales via le standard \textsf{booktabs}.
Ce modèle démontre qu'une publication de haut niveau repose sur une typographie maîtrisée où le gras, l'italique et les couleurs sont employés avec une stricte parcimonie.

\bigskip
\noindent\textbf{Mots-clés :} Lua\LaTeX{}, ClasseGuillaume, typographie scientifique, architecture modulaire, microtypographie, booktabs, siunitx, annotations collaboratives, scrlayer-scrpage.

\bigskip
\bigskip

\selectlanguage{english}
\addchap{Abstract}

This document serves as both an in-depth reference tutorial on modern scientific typesetting with \LaTeX{} and a comprehensive showcase of the custom \texttt{ClasseGuillaume} class.
Its core ambition is to demonstrate how uncompromising academic rigor can seamlessly pair with an elegant, balanced, and airy editorial design.
Every dimension of contemporary academic thesis and engineering report authoring is explored: modular source code architecture, reproducible compilation under the Lua\LaTeX{} engine, intelligent font detection, granular baseline stretch control, and microtypographical penalty optimization.

The reader is guided through practical implementations of the class's exclusive capabilities:
the collaborative marginal review system \texttt{\textbackslash todoH} featuring automatic message truncation and interactive PDF tooltips, on-the-fly page geometry toggling (\texttt{\textbackslash geometrieWord}, \texttt{\textbackslash geometrieLarge}, \texttt{\textbackslash geometrieStandard}), dedicated floating environments for illustrations alongside conventional figures, and publication-ready tables strictly adhering to the \textsf{booktabs} guidelines.
This template illustrates that high-end documents rely on disciplined typesetting where bold weights, italics, and chromatic accents are applied with deliberate parsimony.

\bigskip
\noindent\textbf{Keywords:} Lua\LaTeX{}, ClasseGuillaume, scientific typesetting, modular architecture, microtypography, booktabs, siunitx, collaborative annotations, scrlayer-scrpage.

\selectlanguage{french}


% 5. Avertissement didactique et remerciements
%------------------------------------------------------------------------------------
%  AVERTISSEMENT PÉDAGOGIQUE ET REMERCIEMENTS
%------------------------------------------------------------------------------------
\addchap{Avertissement Didactique}

\begin{boiteSynthese}{Double Nature de ce Document Modèle}
  Ce volume est conçu selon une double perspective indissociable :
  \begin{enumerate}
    \item \textbf{Côté lecteur :} Il se lit de manière fluide comme un manuel d'ingénierie logicielle et de typographie académique contemporaine.
    \item \textbf{Côté auteur / développeur :} Chaque section est un exemple de code source prêt à l'emploi. Le balisage \LaTeX{} des fichiers sous-jacents a été rédigé avec une clarté exemplaire afin de servir de base de copier-coller immédiate pour vos futurs travaux de thèse, d'articles ou de mémoires de fin d'études.
  \end{enumerate}
\end{boiteSynthese}

\vspace{1cm}

\section*{Remerciements Typographiques et Institutionnels}

Ce projet s'appuie sur la richesse de l'écosystème open-source \LaTeX{} et sur les travaux fondateurs de Donald Knuth et Leslie Lamport, perpétués avec brio par les équipes du projet \LaTeX{}, du moteur Lua\LaTeX{} et de la suite KOMA-Script développée par Markus Kohm.

Nos remerciements s'adressent tout particulièrement :
\begin{itemize}
  \item À la communauté du \emph{Comprehensive \TeX{} Archive Network} (CTAN) pour la maintenance pérenne des packages incontournables mis en œuvre dans cette classe (\textsf{fontspec}, \textsf{unicode-math}, \textsf{tcolorbox}, \textsf{biblatex}, \textsf{cleveref}, \textsf{siunitx}).
  \item Aux enseignants, chercheurs et relecteurs de l'UCLouvain dont les retours méthodologiques ont permis d'éprouver les contraintes de mise en page réelles des mémoires de master et thèses de doctorat.
  \item À l'ensemble des contributeurs aux standards de reproductibilité numérique et d'archivage pérenne des connaissances scientifiques.
\end{itemize}



% 6. Table des matières générale
\setcounter{tocdepth}{2}
\tableofcontents

%------------------------------------------------------------------------------------
%  CORPS DU TEXTE (MAINMATTER : Numérotation arabe, activation du filigrane)
%------------------------------------------------------------------------------------
\mainmatter
\afficherLogo % Activation du filigrane pour l'ensemble du corps de document

%====================================================================================
%  PARTIE I : FONDATIONS TYPOGRAPHIQUES, ARCHITECTURE & COULEURS
%====================================================================================
\part{Fondations Typographiques, Architecture \& Couleurs}
\label{part:fondations}

\chapter{Architecture Modulaire et Moteur Lua\LaTeX{}}
\label{chap:fondations}
%------------------------------------------------------------------------------------
%  SECTION 1.1 : ARCHITECTURE MODULAIRE D'UN PROJET LATEX PROFESSIONNEL
%------------------------------------------------------------------------------------
\section{Architecture Modulaire et Séparation des Préoccupations}
\label{sec:architecture_projet}

La rédaction d'un document académique de grande envergure (mémoire, thèse de doctorat, monographie technique) ne peut reposer sur un fichier unique monolithique.
Un projet durable et évolutif obéit au principe cardinal de \textbf{séparation des préoccupations} (\emph{Separation of Concerns}) :
la structure et la typographie sont encapsulées dans la classe de document ou des feuilles de style, tandis que le contenu textuel réside dans des fichiers autonomes et légers.

\subsection{Arborescence standardisée du projet}
\label{subsec:arborescence}

L'arborescence recommandée par ce modèle sépare nettement chaque catégorie d'artefact :

\begin{itemize}
  \item \textbf{\texttt{ClasseGuillaume.cls} :} Le cœur typographique et technique. Il configure les géométries, les en-têtes KOMA-Script, les paquets scientifiques et les environnements exclusifs.
  \item \textbf{\texttt{main.tex} :} Le chef d'orchestre. Ce fichier ne contient presque aucun texte direct, mais uniquement la déclaration de la classe, les métadonnées globales, le chargement des annexes et la séquence d'inclusion des chapitres via la commande \texttt{\textbackslash input\{...\}}.
  \item \textbf{\texttt{chapters/} :} Les sources textuelles organisées par répertoires thématiques numérotés (\texttt{00\_frontmatter/}, \texttt{01\_fondations/}, etc.).
  \item \textbf{\texttt{tables/} :} Les tableaux complexes ou volumineux isolés dans des fichiers autonomes pour maintenir une parfaite lisibilité du texte principal.
  \item \textbf{\texttt{bib/references.bib} :} Le corpus bibliographique centralisé au format Bib\LaTeX{}.
  \item \textbf{\texttt{images/} :} Les illustrations vectorielles (PDF/SVG), schémas techniques et photographies matricielles.
  \item \textbf{\texttt{scripts/} :} Les scripts d'automatisation de compilation (\texttt{compile.ps1}) et de nettoyage des fichiers temporaires (\texttt{clean.ps1}).
  \item \textbf{\texttt{build/} :} Le répertoire dédié à l'isolation de tous les fichiers auxiliaires générés par le compilateur (\texttt{.aux}, \texttt{.bcf}, \texttt{.log}, \texttt{.toc}).
\end{itemize}

\subsection{Le paradigme Git-friendly : Semantic Line Breaks}
\label{subsec:semantic_breaks}

Dans le développement logiciel moderne, la gestion de versions avec Git est indispensable.
Cependant, un document \LaTeX{} dont les paragraphes sont rédigés en un seul bloc kilométrique rend les comparaisons de versions (\texttt{git diff}) totalement illisibles : la moindre correction d'une virgule marque l'ensemble du paragraphe comme modifié.

Pour remédier à cet écueil, ce modèle applique rigoureusement la convention des \textbf{sauts de ligne sémantiques} (\emph{semantic line breaks}) :
\begin{boiteSynthese}{Règle d'or de rédaction Git-friendly}
  Chaque phrase complète ou proposition grammaticale majeure doit débuter sur une nouvelle ligne dans le fichier source \texttt{.tex}.
  Pour le compilateur \LaTeX{}, un simple retour à la ligne est assimilé à une espace ordinaire entre les mots ; le saut de paragraphe ne s'effectue qu'après une ligne entièrement vierge.
  Cette pratique assure des revues de code claires, des fusions de branches sans conflit et un suivi historique précis de chaque argument rédigé.
\end{boiteSynthese}

\subsection{Compilation automatisée et hygiène du répertoire source}
\label{subsec:compilation_build}

La compilation manuelle enchaînant plusieurs passes (\texttt{lualatex}, puis \texttt{biber}, puis deux nouvelles passes de \texttt{lualatex}) est fastidieuse et source d'erreurs d'incohérence dans les références croisées.

La classe s'appuie sur l'outil de référence \texttt{latexmk}, configuré via le fichier masqué \texttt{.latexmkrc} à la racine du projet :
\begin{enumerate}
  \item L'option \texttt{-outdir=build} redirige la totalité des artefacts intermédiaires dans le sous-dossier \texttt{build/}.
  \item Le répertoire racine reste ainsi vierge de toute pollution par des fichiers parasites (\texttt{.aux}, \texttt{.log}, \texttt{.fls}).
  \item \texttt{latexmk} détecte automatiquement les modifications apportées aux dépendances (bibliographies, graphiques, fichiers inclus) et calcule exactement le nombre optimal de passes nécessaires pour stabiliser le document.
\end{enumerate}


%------------------------------------------------------------------------------------
%  SECTION 1.2 : LE MOTEUR LUALATEX ET LA GESTION MODERNE DES POLICES
%------------------------------------------------------------------------------------
\section{Moteur Lua\LaTeX{} et Gestion Dynamique des Polices}
\label{sec:lualatex_polices}

Pendant plusieurs décennies, le compilateur historique \textsf{pdf\LaTeX{}} a dominé le paysage éditorial.
Cependant, ses limitations intrinsèques liées à l'encodage sur 8 bits nécessitaient une armada de paquets intermédiaires (\texttt{inputenc}, \texttt{fontenc}, \texttt{textcomp}) et rendaient l'utilisation des polices vectorielles modernes OpenType particulièrement laborieuse.

Le modèle \texttt{ClasseGuillaume} s'appuie résolument et exclusivement sur \textbf{Lua\LaTeX{}}, successeur moderne alliant un moteur \TeX{} natif Unicode à la puissance du langage de script embarqué Lua.

\subsection{L'abandon des paquets d'encodage obsolètes}
\label{subsec:unicode_natif}

Sous Lua\LaTeX{}, le code source \LaTeX{} est traité nativement en \textbf{UTF-8}.
Les instructions suivantes, jadis ubiquitaires, sont devenues \textbf{strictement obsolètes} et ne doivent jamais figurer dans votre document :
\begin{itemize}
  \item \texttt{\textbackslash usepackage[utf8]\{inputenc\}} : Inutile et néfaste, Lua\LaTeX{} gère l'Unicode dès sa première ligne.
  \item \texttt{\textbackslash usepackage[T1]\{fontenc\}} : Obsolète, le système de polices \textsf{fontspec} régit directement le rendu glyphe par glyphe.
\end{itemize}

Grâce à ce traitement natif, tous les caractères accentués français (é, è, à, ç, œ, « ») ainsi que les symboles scientifiques et alphabets internationaux sont insérés sans aucune macro d'échappement.

\subsection{Les trois modes de chargement des polices typographiques}
\label{subsec:trois_modes_polices}

Une question cruciale lors de la transition entre une plateforme cloud comme Overleaf et une installation locale (TeX Live) concerne l'accessibilité des polices de caractères :

\begin{enumerate}
  \item \textbf{Mode 1 : Les polices système OS (la méthode la plus directe) :}
  Sur une machine équipée de Windows ou macOS, le sous-système \textsf{luaotfload} indexe automatiquement le dossier des polices système (\texttt{C:\textbackslash Windows\textbackslash Fonts\textbackslash}).
  N'importe quelle police installée sur votre ordinateur (telle que \textsf{Calibri}, \textsf{Segoe UI}, \textsf{Arial}, \textsf{Consolas}) peut être appelée directement par son nom usuel :
  \begin{verbatim}
  \setmainfont{Calibri}
  \setsansfont{Calibri}
  \end{verbatim}
  Le compilateur associe automatiquement les variantes grasses (\texttt{calibrib.ttf}), italiques (\texttt{calibrii.ttf}) et grasses-italiques sans aucune intervention manuelle.
  Aucun fichier de police n'a besoin d'être copié dans le projet !

  \item \textbf{Mode 2 : Les polices OpenType natives de \TeX{} Live (universelles et libres) :}
  La distribution complète \TeX{} Live installe sur votre machine des centaines de familles de polices professionnelles prêtes à l'emploi.
  Par exemple, la suite \textsf{\TeX{} Gyre} fournit des équivalents modernes et fidèles des plus grands classiques :
  \begin{itemize}
    \item \textsf{TeX Gyre Heros} (équivalent moderne d'Helvetica),
    \item \textsf{TeX Gyre Termes} (équivalent de Times),
    \item \textsf{TeX Gyre Pagella} (équivalent de Palatino),
    \item \textsf{Libertinus Serif} et \textsf{Fira Sans}.
  \end{itemize}
  Ces polices fonctionnent de manière identique et immédiate sur Windows, Linux, macOS et Overleaf, sans aucun fichier physique dans le répertoire du projet.

  \item \textbf{Mode 3 : Le repli local dans le projet (\texttt{./Fonts/...}) (méthode Overleaf) :}
  Sur les serveurs cloud comme Overleaf (qui tournent sous des conteneurs Linux), les polices propriétaires Microsoft ne sont pas préinstallées sous licence libre.
  Pour pallier cette absence sans faire échouer la compilation, \texttt{ClasseGuillaume} embarque un dossier local de secours \texttt{./Fonts/}.
  Grâce à la commande conditionnelle intelligente \texttt{\textbackslash IfFontExistsTF}, la classe interroge d'abord le système de votre PC ; si la police y est absente (ex. sur Overleaf), elle active automatiquement le fichier de repli local.
\end{enumerate}

\subsection{Options de classe disponibles pour le choix de la police}
\label{subsec:options_polices}

Dans l'en-tête de votre document principal \texttt{main.tex}, vous pouvez sélectionner votre famille de polices en un mot-clé au sein de la commande \texttt{\textbackslash documentclass} :

\begin{itemize}
  \item \texttt{font=calibri} *(option par défaut)* : Active \textsf{Calibri} pour un style sobre, fluide et moderne, hautement lisible sur écran comme à l'impression.
  \item \texttt{font=heros} : Active \textsf{TeX Gyre Heros}, conférant un rendu suisse net et contemporain.
  \item \texttt{font=libertinus} : Active \textsf{Libertinus Serif}, idéale pour les mémoires littéraires ou les thèses académiques traditionnelles à fort contenu textuel.
  \item \texttt{font=comicsans} : Active \textsf{Comic Sans MS} (utile à des fins pédagogiques ou pour des présentations informelles).
  \item \texttt{font=standard} : Active la police standard \textsf{Latin Modern}.
\end{itemize}

\subsection{Démonstration des commandes de polices locales}
\label{subsec:commandes_polices_locales}

Outre la police générale du document, la classe met à votre disposition des commandes spécifiques pour des encarts, titres ou spécimens ponctuels :

\begin{itemize}
  \item \textbf{Police de titre de chapitre :} La commande \texttt{\textbackslash AsapFont} applique la police {\AsapFont Asap Condensed}, conférant un impact visuel fort aux intitulés de chapitres.
  \item \textbf{Spécimen Calibri :} Accessible ponctuellement via {\CalibriFont \textbackslash CalibriFont pour forcer le style sans-serif}.
  \item \textbf{Spécimen fantaisie :} Accessible via {\ComicSansMSFont \textbackslash ComicSansMSFont pour des démonstrations ciblées}.
\end{itemize}


%------------------------------------------------------------------------------------
%  SECTION 1.3 : GESTION DES INTERLIGNES ET MICROTYPOGRAPHIE
%------------------------------------------------------------------------------------
\section{Interlignage Contrôlé et Microtypographie Avancée}
\label{sec:interlignes_microtype}

La lisibilité d'un document long repose sur une harmonie subtile entre la taille de la police, la longueur de la ligne et l'espacement vertical entre deux lignes consécutives.
Trop serré, le texte fatigue l'œil lors des retours à la ligne ; trop espacé, la page semble déstructurée et les paragraphes perdent leur cohésion.

\subsection{Le défi classique de l'interligne 1,5 et la réponse de la classe}
\label{subsec:defi_interligne}

La majorité des universités et jurys de thèse exigent un interligne de 1,5 pour le confort de lecture et l'annotation manuscrite des épreuves.
Le paquet standard \textsf{setspace} fournit la commande \texttt{\textbackslash onehalfspacing}.

Cependant, en \LaTeX{} standard (\emph{vanilla}), l'activation globale de cet interligne engendre un effet visuel désastreux :
lorsqu'un titre de section ou de chapitre s'étend sur deux lignes, l'interligne 1,5 lui est également appliqué !
Le titre apparaît alors artificiellement distendu et déconnecté de son texte.

\begin{boiteSynthese}{Protection automatique des titres dans ClasseGuillaume}
  Pour pallier cette anomalie, \texttt{ClasseGuillaume} intègre une règle typographique stricte au sein des polices KOMA-Script :
  \begin{verbatim}
  \addtokomafont{section}{\setstretch{1}\selectfont}
  \addtokomafont{subsection}{\setstretch{1}\selectfont}
  \addtokomafont{subsubsection}{\setstretch{1}\selectfont}
  \end{verbatim}
  Ainsi, le corps du mémoire bénéficie de la fluidité de l'interligne 1,5, tandis que l'ensemble des titres longs conservent un interligne simple compact, rigoureux et professionnel.
\end{boiteSynthese}

\subsection{Microtypographie et élimination des lignes orphelines et veuves}
\label{subsec:microtypographie}

Un document scientifique soigné ne doit jamais laisser une ligne isolée au bas d'une page (ligne orpheline) ou en haut de la page suivante (ligne veuve).
De même, une césure coupant un mot à la toute dernière ligne d'une page est considérée comme un manquement aux règles de l'art typographique.

Le paquet \textsf{microtype} associé à la configuration interne de \texttt{ClasseGuillaume} résout ces écueils par un ajustement microscopique de l'espacement inter-mots et par des pénalités strictes :

\begin{itemize}
  \item \texttt{\textbackslash clubpenalty = 10000} : Interdit formellement la présence d'une ligne orpheline en bas de page.
  \item \texttt{\textbackslash widowpenalty = 10000} : Interdit le rejet d'une ligne veuve en haut de page.
  \item \texttt{\textbackslash displaywidowpenalty = 10000} : Protège l'agencement des formules mathématiques détachées.
  \item \texttt{\textbackslash brokenpenalty = 10000} : Empêche la césure d'un mot à cheval sur deux pages consécutives.
  \item \texttt{\textbackslash interfootnotelinepenalty = 9999} : Empêche qu'une note de bas de page ne soit scindée arbitrairement sur deux pages.
\end{itemize}

Grâce à ces paramètres, le compilateur Lua\LaTeX{} recalcule subtilement les blancs typographiques de chaque paragraphe pour produire des pages équilibrées sans rupture visuelle.



\chapter{Hiérarchie des Titres, Typographie Textuelle \& Charte Chromatique}
\label{chap:titres_couleurs}
%------------------------------------------------------------------------------------
%  SECTION 2.1 : DÉMONSTRATION DE TOUS LES NIVEAUX DE SECTIONNEMENT
%------------------------------------------------------------------------------------
\section{Hiérarchie Intégrale des Titres et Divisions Structurelles}
\label{sec:niveaux_sectionnement}

Un document scientifique clair et pédagogique doit s'appuyer sur une structure hiérarchique sans ambiguïté.
La classe \texttt{ClasseGuillaume}, construite sur le socle KOMA-Script (\texttt{scrbook}), personnalise l'ensemble des niveaux de sectionnement du document, depuis la grande division générale jusqu'au paragraphe en ligne.

\subsection{Le niveau macroscopique : Les Parties (\texttt{\textbackslash part})}
\label{subsec:parties}

La commande \texttt{\textbackslash part\{Titre de la Partie\}} regroupe plusieurs chapitres au sein d'une unité thématique majeure.
Dans ce modèle, chaque partie génère une page de garde majestueuse, centrée et aérée, portant un numéro romain.
Les parties permettent d'architecturer les documents de plus de 60 pages sans saturer la numérotation des chapitres.

\subsection{Le niveau capitulaire : Les Chapitres numérotés et non numérotés}
\label{subsec:chapitres}

Les chapitres constituent le cœur du découpage :
\begin{itemize}
  \item \textbf{Chapitre numéroté (\texttt{\textbackslash chapter\{...\}}) :}
  La classe personnalise le titre en adoptant la police {\AsapFont Asap Condensed}, la couleur primaire du document, et une numérotation en chiffres romains (\textsc{Chapitre I}, \textsc{Chapitre II}).
  Le titre est harmonieusement centré et séparé du texte par des espacements verticaux soigneusement calibrés (\texttt{beforeskip = -10pt}, \texttt{afterskip = 20pt}).
  \item \textbf{Chapitre non numéroté KOMA-Script (\texttt{\textbackslash addchap\{...\}}) :}
  En \LaTeX{} standard, la commande \texttt{\textbackslash chapter*\{...\}} présente un défaut majeur : elle n'insère pas le chapitre dans la table des matières et ne met pas à jour les en-têtes de page.
  La commande KOMA-Script \texttt{\textbackslash addchap\{...\}} résout parfaitement ce problème : elle produit un chapitre sans numéro, mais l'enregistre automatiquement dans la table des matières et actualise les en-têtes (tel que réalisé pour le Résumé et les Remerciements).
\end{itemize}

\subsection{Les divisions intermédiaires : Sections et Sous-sections}
\label{subsec:sections_subsections}

La numérotation des sections adopte une convention hiérarchique arabe préfixée par le numéro du chapitre courant :
\begin{itemize}
  \item \textbf{Section (\texttt{\textbackslash section\{...\}}) :}
  Numérotée par exemple \textbf{II.1}, mise en valeur par la couleur secondaire (\texttt{secondaryColor}) et un corps de texte large en gras.
  \item \textbf{Section non numérotée dans la TOC (\texttt{\textbackslash addsec\{...\}}) :}
  À l'instar de \texttt{\textbackslash addchap}, elle insère un en-tête de section sans numéro tout en figurant dans la table des matières.
  \item \textbf{Sous-section (\texttt{\textbackslash subsection\{...\}}) :}
  Numérotée par exemple \textbf{II.1.1}, typographiée en gris foncé soutenu (\texttt{darkgray}) pour créer un contraste hiérarchique doux.
  \item \textbf{Sous-sous-section (\texttt{\textbackslash subsubsection\{...\}}) :}
  Offre un niveau de granularité supplémentaire pour les démonstrations techniques pointues ou les descriptions expérimentales détaillées.
\end{itemize}

Voici ci-dessous une démonstration immédiate d'une sous-sous-section :

\subsubsection{Exemple pratique de sous-sous-section technique}
Cette sous-sous-section permet de délimiter un point précis d'analyse sans surcharger la numérotation globale.
Elle s'intègre harmonieusement dans le flux de lecture sans créer de rupture visuelle excessive.

\subsection{Les divisions terminales : Paragraphes avec puce en ligne}
\label{subsec:paragraphes_lignes}

En \LaTeX{} standard, l'environnement \texttt{\textbackslash paragraph\{...\}} engendre fréquemment une mise en page gauche : soit il renvoie maladroitement à la ligne, soit il s'accole de manière peu esthétique.

\begin{boiteSynthese}{L'innovation de ClasseGuillaume pour \texttt{\textbackslash paragraph}}
  La classe redéfinit entièrement la commande \texttt{\textbackslash paragraph} :
  \begin{verbatim}
  \RedeclareSectionCommand[
    indent = 5mm,
    afterskip = -1em,
    font = {\color{primaryColor}\bfseries\textbullet~}
  ]{paragraph}
  \end{verbatim}
  Le titre du paragraphe est automatiquement précédé d'une \textbf{puce discrète colorée} (\texttt{\textbullet}), indenté de \qty{5}{\milli\meter}, et le texte du corps s'enchaîne \textbf{directement sur la même ligne} après une espace fine.
\end{boiteSynthese}

Démonstration en situation réelle :

\paragraph{Critère de modularité}
Le code source est segmenté en unités logiques autonomes pour autoriser une maintenance aisée et des révisions collaboratives sans friction.

\paragraph{Critère d'accessibilité typographique}
Les contrastes de couleurs, la lisibilité des polices et l'aération spatiale garantissent un confort de lecture optimal pour les jurys et relecteurs.

\subparagraph{Niveau sub-paragraph}
Pour les nomenclatures ou sous-critères nécessitant une hiérarchisation ultime, la commande \texttt{\textbackslash subparagraph} offre une indentation spécifique tout en conservant une discrétion absolue.


%------------------------------------------------------------------------------------
%  SECTION 2.2 : TYPOGRAPHIE TEXTUELLE ET RÈGLE DE PARCIMONIE
%------------------------------------------------------------------------------------
\section{Typographie Textuelle et Règle d'Or de la Parcimonie}
\label{sec:typographie_parcimonie}

L'une des erreurs les plus fréquentes dans la rédaction technique consiste à abuser des variations de graisse (gras) et de style (italique, soulignement, surlignage).
Comme le rappelait Robert Bringhurst dans \emph{The Elements of Typographic Style}~\cite{bringhurst2004elements}, \textbf{« quand tout est accentué, plus rien n'est visible »}.
La crédibilité d'un document académique repose sur une sobriété rigoureuse et une sémantique précise.

\subsection{L'usage raisonné du gras, de l'italique et des petites capitales}
\label{subsec:usages_styles}

Chaque modification typographique doit répondre à un rôle sémantique précis :

\begin{itemize}
  \item \textbf{Le gras (\texttt{\textbackslash textbf\{...\}}) :}
  Il doit être réservé exclusivement aux \textbf{mots-clés fondamentaux} et aux concepts pivots lors de leur première introduction.
  Une phrase entière ne doit jamais être mise en gras, sous peine de détruire le gris typographique de la page.
  
  \item \textbf{L'emphase sémantique (\texttt{\textbackslash emph\{...\}}) :}
  Elle signale une inflexion de la pensée dans le discours.
  L'avantage majeur de \texttt{\textbackslash emph} est qu'elle s'inverse automatiquement : dans un passage déjà en italique, un mot englobé dans \texttt{\textbackslash emph} repasse en droit (romain), préservant le contraste.

  \item \textbf{L'italique factuel (\texttt{\textbackslash textit\{...\}}) :}
  Il s'applique aux locutions latines d'usage scientifique (\textit{a priori}, \textit{a posteriori}, \textit{in situ}, \textit{in vivo}, \textit{de facto}), aux termes empruntés à une langue étrangère non francisée, ainsi qu'aux titres d'ouvrages cités dans le texte.

  \item \textbf{La chasse fixe (\texttt{\textbackslash texttt\{...\}}) :}
  Réservée aux commandes \LaTeX{}, noms de packages (\texttt{biblatex}, \texttt{siunitx}), chemins de fichiers et variables de code.

  \item \textbf{Les petites capitales (\texttt{\textbackslash textsc\{...\}}) :}
  Idéales pour les patronymes d'auteurs (ex. Guillaume \textsc{Winant}) ou certains sigles académiques, évitant que des majuscules massives n'alourdissent le paragraphe.
\end{itemize}

\subsection{Le surlignage sécurisé sous Lua\LaTeX{} avec \textsf{lua-ul}}
\label{subsec:surlignage_luaul}

En \LaTeX{} classique, le soulignement et le surlignage reposaient sur des paquets comme \texttt{soul} ou \texttt{ulem}.
Ces paquets sont tristement célèbres pour leurs plantages systématiques dès qu'un mot contient un caractère accentué UTF-8, une césure ou une référence croisée.

\texttt{ClasseGuillaume} résout définitivement ce problème en intégrant le paquet moderne \textsf{lua-ul} couplé à \textsf{luacolor} :
\begin{boiteSynthese}{Surlignage non destructif avec \texttt{\textbackslash highLight}}
  La commande \texttt{\textbackslash highLight[couleur]\{texte\}} permet de surligner délicatement un passage clé.
  Elle gère parfaitement les sauts de ligne, les césures de mots et tous les caractères Unicode :
  Par exemple, voici un passage \highLight[yellow!35]{surligné en jaune pastel avec césure automatique}, ou encore \highLight[secondaryColor!20]{un surlignage teinté aux couleurs de la charte graphique}.
\end{boiteSynthese}
Ce mécanisme forme la fondation technique de la commande collaborative \texttt{\textbackslash todoH} présentée au Chapitre~\ref{chap:collaboration_todos}.


%------------------------------------------------------------------------------------
%  SECTION 2.3 : CHARTE GRAPHIQUE ET THÉORIE DES COULEURS
%------------------------------------------------------------------------------------
\section{Charte Graphique, Théorie des Couleurs et Personnalisation}
\label{sec:charte_chromatique}

Une publication scientifique contemporaine ne se contente plus du noir et blanc austère.
Une utilisation ciblée de la couleur renforce la hiérarchie visuelle, guide la navigation du lecteur et met en valeur les résultats majeurs.
Toutefois, la palette doit être rigoureusement structurée pour éviter tout effet « sapin de Noël » incompatible avec le sérieux universitaire.

\subsection{La palette chromatique par défaut de \texttt{ClasseGuillaume}}
\label{subsec:palette_defaut}

La classe s'appuie sur le paquet \textsf{xcolor} (avec les modèles de noms \texttt{dvipsnames}) et définit quatre couleurs maîtresses :

\begin{itemize}
  \item \textbf{\texttt{primaryColor} (\textcolor{NavyBlue}{\rule{8pt}{8pt}} \texttt{NavyBlue}) :}
  Couleur dominante du document. Elle habille les grands titres de chapitres, les encadrements de boîtes principales, la puce des paragraphes et les liens de navigation interne.
  
  \item \textbf{\texttt{secondaryColor} (\textcolor{TealBlue}{\rule{8pt}{8pt}} \texttt{TealBlue}) :}
  Couleur secondaire complémentaire. Elle structure les titres de sections et sous-titres, apportant une respiration visuelle plus douce.

  \item \textbf{\texttt{accentColor} (\textcolor{Bittersweet}{\rule{8pt}{8pt}} \texttt{Bittersweet}) :}
  Couleur d'accentuation chaude. Utilisée avec modération pour signaler des alertes méthodologiques ou des points de vigilance.

  \item \textbf{\texttt{mdtRed} (\textcolor{red!50!black}{\rule{8pt}{8pt}} \texttt{red!50!black}) :}
  Rouge bordeaux assombri, réservé exclusivement aux liens URL externes dans le texte et la bibliographie, évitant le rouge vif agressif par défaut d'\textsf{hyperref}.
\end{itemize}

\subsection{Comment personnaliser l'ambiance du document en 4 lignes}
\label{subsec:personnalisation_couleurs}

Grâce à l'usage de la macro \texttt{\textbackslash colorlet}, vous n'avez jamais besoin de modifier le code interne de la classe pour adapter le mémoire à la charte de votre université ou entreprise.
Il vous suffit d'insérer dans votre préambule (dans \texttt{main.tex}) :

\begin{verbatim}
% Exemple : Bascule vers une palette 'Bordeaux & Ardoise'
\colorlet{primaryColor}{Maroon}
\colorlet{secondaryColor}{SlateGrey}
\colorlet{accentColor}{BurntOrange}
\colorlet{mdtRed}{Maroon!80!black}
\end{verbatim}

\subsection{Démonstration visuelle de palettes académiques harmonieuses}
\label{subsec:nuanciers_harmonieux}

Voici un aperçu comparatif de combinaisons chromatiques professionnelles testées et recommandées :

\begin{table}[htbp]
  \centering
  \small
  \caption{Palettes recommandées selon le domaine scientifique}
  \label{tab:palettes_recommandees}
  \begin{tabularx}{\textwidth}{l p{3.5cm} p{4cm} X}
    \toprule
    \textbf{Nom de la Palette} & \textbf{Couleur Primaire} & \textbf{Couleur Secondaire} & \textbf{Domaine Idéal} \\
    \midrule
    \textbf{Océan Académique} (Défaut) & \textcolor{NavyBlue}{\rule{6pt}{6pt}} Bleu Marine (\texttt{NavyBlue}) & \textcolor{TealBlue}{\rule{6pt}{6pt}} Sarcelle (\texttt{TealBlue}) & Médecine, Droit, Économie \\
    \textbf{Prestige Humanités} & \textcolor{Maroon}{\rule{6pt}{6pt}} Bordeaux (\texttt{Maroon}) & \textcolor{darkgray}{\rule{6pt}{6pt}} Anthracite (\texttt{darkgray}) & Histoire, Philosophie, Lettres \\
    \textbf{Sciences Vivant} & \textcolor{ForestGreen}{\rule{6pt}{6pt}} Vert Forêt (\texttt{ForestGreen}) & \textcolor{CadetBlue}{\rule{6pt}{6pt}} Bleu Cadre (\texttt{CadetBlue}) & Biologie, Agronomie, Écologie \\
    \textbf{Tech \& Ingénierie} & \textcolor{RoyalBlue}{\rule{6pt}{6pt}} Bleu Roi (\texttt{RoyalBlue}) & \textcolor{Bittersweet}{\rule{6pt}{6pt}} Ambre (\texttt{Bittersweet}) & Informatique, Physique, Aéro \\
    \bottomrule
  \end{tabularx}
\end{table}

Chacune de ces palettes garantit un ratio de contraste supérieur aux normes d'accessibilité visuelle WCAG AA, que ce soit pour une lecture sur écran haute résolution ou lors d'une impression sur papier couché.
Un guide complet de personnalisation avec code source est fourni en Annexe~\ref{annexe:guide_palettes}.



%====================================================================================
%  PARTIE II : SCIENCES, MATHÉMATIQUES & TABLEAUX PROFESSIONNELS
%====================================================================================
\part{Sciences, Mathématiques \& Tableaux Professionnels}
\label{part:sciences_tables}

\chapter{Sciences, Typographie Mathématique et Unités Internationales}
\label{chap:sciences_maths}
%------------------------------------------------------------------------------------
%  SECTION 3.1 : FORMULES MATHÉMATIQUES MODERNES AVEC UNICODE-MATH
%------------------------------------------------------------------------------------
\section{Typographie Mathématique Moderne et \textsf{unicode-math}}
\label{sec:formules_mathematiques}

La composition mathématique constitue historiquement le point fort absolu de \LaTeX{}.
Cependant, de nombreuses pratiques héritées des années 1980 persistent encore à tort dans les documents contemporains.
\texttt{ClasseGuillaume} intègre nativement le paquet \textsf{unicode-math} couplé à la police vectorielle mathématique de référence \textbf{XITS Math}, assurant un rendu irréprochable des symboles et une harmonisation d'échelle avec le texte environnant (\texttt{Scale=MatchLowercase}).

\subsection{L'interdiction absolue de \texttt{\$\$...\$\$} et les environnements recommandés}
\label{subsec:interdiction_dollars}

L'une des règles fondamentales de \LaTeXe{} moderne concerne les formules hors-texte :
\begin{boiteSynthese}{Pourquoi la syntaxe \texttt{\$\$...\$\$} est proscrite}
  La syntaxe \texttt{\$\$...\$\$} est une relique primitive du moteur \TeX{} originel (1978).
  En \LaTeX{}, elle perturbe gravement le calcul vertical des marges (\texttt{abovedisplayskip}), court-circuite les options de classes telles que \texttt{fleqn}, et interdit toute numérotation automatique ou référence croisée.
  
  \textbf{Règles à appliquer sans exception :}
  \begin{itemize}
    \item Pour une équation hors-texte non numérotée : utilisez \texttt{\textbackslash [ ... \textbackslash ]}.
    \item Pour une équation hors-texte numérotée : utilisez \texttt{\textbackslash begin\{equation\} ... \textbackslash end\{equation\}}.
    \item Pour des développements multi-lignes alignés : utilisez \texttt{align} ou \texttt{align*}.
  \end{itemize}
\end{boiteSynthese}

\subsection{Ponctuation et texte au sein des équations}
\label{subsec:ponctuation_maths}

Une équation mathématique ne flotte pas dans le vide : elle constitue une phrase grammaticale à part entière ou s'insère dans une proposition existante.
Par conséquent, elle doit impérativement comporter sa \textbf{ponctuation finale} (virgule ou point), placée à l'intérieur de l'environnement mathématique.

De surcroît, les indices ou mots insérés dans une formule ne doivent jamais être saisis en mode mathématique pur, sous peine d'être interprétés comme le produit de variables juxtaposées ($vitesse$ au lieu de $\text{vitesse}$).
On utilise systématiquement la commande \texttt{\textbackslash text\{...\}} du paquet \textsf{mathtools} :

\begin{equation}
  \label{eq:bernoulli_moderne}
  p_1 + \frac{1}{2} \rho v_1^2 + \rho g z_1 = p_2 + \frac{1}{2} \rho v_2^2 + \rho g z_2 + \Delta p_{\text{pertes}},
\end{equation}
où $p$ désigne la pression statique, $\rho$ la masse volumique du fluide, et $\Delta p_{\text{pertes}}$ les pertes de charge singulières et régulières.

\subsection{Démonstration de structures mathématiques complexes}
\label{subsec:structures_complexes}

L'environnement \texttt{align} permet de coordonner des démonstrations pas à pas grâce au point d'alignement \texttt{\&} :
\begin{align}
  \nabla \times \mathbf{E} &= -\frac{\partial \mathbf{B}}{\partial t}, \label{eq:maxwell_faraday} \\
  \nabla \times \mathbf{B} &= \mu_0 \mathbf{J} + \mu_0 \varepsilon_0 \frac{\partial \mathbf{E}}{\partial t}. \label{eq:maxwell_ampere}
\end{align}
Une référence croisée intelligente via \textsf{cleveref} permet d'appeler simultanément les~\cref{eq:maxwell_faraday,eq:maxwell_ampere} sans répétition lourde.

Pour les systèmes différentiels avec conditions aux limites ou fonctions définies par morceaux, l'environnement \texttt{cases} de \textsf{amsmath} garantit un rendu net :
\begin{equation}
  \label{eq:fonction_morceaux}
  f(x, t) = 
  \begin{cases}
    \int_0^t \mathrm{e}^{-\lambda \tau} \cos(\omega \tau) \, \mathrm{d}\tau & \text{si } x \ge 0, \\[2mm]
    \sum_{n=1}^{\infty} \frac{(-1)^n}{n^2 + 1} \sin\left(\frac{n \pi x}{L}\right) & \text{si } x < 0.
  \end{cases}
\end{equation}
Remarquez l'écriture droite du symbole différentiel $\mathrm{d}\tau$ et de la constante d'Euler $\mathrm{e}$, conformément aux normes internationales ISO 80000-2.

Les calculs matriciels s'expriment avec clarté via \texttt{pmatrix} ou \texttt{bmatrix} :
\begin{equation}
  \label{eq:matrice_rotation}
  \mathbf{R}_z(\theta) = 
  \begin{bmatrix}
    \cos\theta & -\sin\theta & 0 \\
    \sin\theta &  \cos\theta & 0 \\
    0          &  0          & 1
  \end{bmatrix},
  \quad
  \det\left(\mathbf{R}_z(\theta)\right) = 1.
\end{equation}
L'\cref{eq:matrice_rotation} illustre la régularité typographique des délimiteurs étirables automatiques.


%------------------------------------------------------------------------------------
%  SECTION 3.2 : GESTION DES UNITÉS ET GRANDEURS PHYSIQUES AVEC SIUNITX
%------------------------------------------------------------------------------------
\section{Réglementation Internationale des Unités et \textsf{siunitx}}
\label{sec:unites_siunitx}

L'écriture des grandeurs physiques et des unités de mesure est régie au niveau mondial par le Bureau International des Poids et Mesures (BIPM) et la norme ISO 80000-1.
Dans un rapport universitaire d'ingénierie ou de médecine, saisir manuellement les unités (ex. \texttt{25 kg}) constitue une faute technique majeure :
\begin{itemize}
  \item Une coupure de ligne accidentelle peut séparer la valeur numérique de son symbole (ex. « 25 » au bout de la ligne et « kg » sur la ligne suivante).
  \item En mode mathématique (\texttt{\$25 km/h\$}), les lettres sont composées en italique mathématique ($km/h$), symbolisant le produit $k \times m / h$ et non l'unité kilomètre par heure.
  \item La séparation des milliers et la virgule décimale varient selon la langue du document.
\end{itemize}

Le paquet \textsf{siunitx} résout l'intégralité de ces problématiques en automatisant le respect strict des conventions typographiques.

\subsection{Configuration francophone native de \texttt{ClasseGuillaume}}
\label{subsec:config_siunitx}

La classe intègre une configuration francophone rigoureuse dans son préambule :
\begin{verbatim}
\sisetup{
  locale = FR,                % Virgule décimale et espace fin entre milliers
  mode = match,               % Harmonise la police avec le texte environnant
  per-mode = symbol,          % Affiche m/s plutôt que m.s^-1
  range-phrase = { à },       % Plage lisible : 10 à 20 km
  range-units = single        % Évite la répétition lourde : 10 à 20 km
}
\end{verbatim}

\subsection{Les commandes indispensables en situation réelle}
\label{subsec:commandes_siunitx}

Le tableau synthétique suivant récapitule les quatre commandes maîtresses du paquet :

\begin{table}[htbp]
  \centering
  \small
  \caption{Commandes fondamentales du paquet \textsf{siunitx}}
  \label{tab:commandes_siunitx}
  \begin{tabularx}{\textwidth}{p{2.6cm} p{6.2cm} X}
    \toprule
    \textbf{Commande \LaTeX{}} & \textbf{Code source type} & \textbf{Rendu typographique obtenu} \\
    \midrule
    \texttt{\textbackslash qty} & \texttt{\textbackslash qty\{250.5\}\{\textbackslash kilo\textbackslash pascal\}} & Une pression de \qty{250.5}{\kilo\pascal} \\
    \texttt{\textbackslash qty} & \texttt{\textbackslash qty\{37.2\}\{\textbackslash degreeCelsius\}} & Une température de \qty{37.2}{\degreeCelsius} \\
    \texttt{\textbackslash qty} & \texttt{\textbackslash qty\{4.7\}\{\textbackslash kilo\textbackslash ohm\}} & Une résistance de \qty{4.7}{\kilo\ohm} \\
    \texttt{\textbackslash qtyrange} & \texttt{\textbackslash qtyrange\{10\}\{25\}\{\textbackslash milli\textbackslash meter\}} & Un diamètre variant de \qtyrange{10}{25}{\milli\meter} \\
    \texttt{\textbackslash unit} & \texttt{\textbackslash unit\{\textbackslash joule\textbackslash per\textbackslash mole\textbackslash per\textbackslash kelvin\}} & L'entropie molaire en \unit{\joule\per\mole\per\kelvin} \\
    \texttt{\textbackslash num} & \texttt{\textbackslash num\{1234567.89\}} & Le nombre formaté \num{1234567.89} \\
    \texttt{\textbackslash num} & \texttt{\textbackslash num\{6.02214076e23\}} & La constante d'Avogadro $N_A = \num{6.02214076e23}$ \\
    \bottomrule
  \end{tabularx}
\end{table}

\subsection{Gestion des angles et pourcentages}
\label{subsec:angles_pourcentages}

La mesure des angles géométriques et des proportions s'effectue avec la même rigueur :
\begin{itemize}
  \item Un angle de déclinaison : \texttt{\textbackslash ang\{45;30;15\}} produit un angle de \ang{45;30;15}.
  \item Un pourcentage sans ambiguïté : \texttt{\textbackslash qty\{98.5\}\{\textbackslash percent\}} donne \qty{98.5}{\percent}, assurant l'espace fine insécable obligatoire en typographie française entre le nombre et le symbole pour cent.
\end{itemize}
Ces mécanismes évitent tout risque d'incohérence entre les différents chapitres d'un projet rédigé en équipe.



\chapter{Tableaux Scientifiques et Typographie Tabulaire}
\label{chap:tableaux}
%------------------------------------------------------------------------------------
%  SECTION 4.1 : LES RÈGLES D'OR DE LA TYPOGRAPHIE TABULAIRE AVEC BOOKTABS
%------------------------------------------------------------------------------------
\section{Les Règles d'Or de la Typographie Tabulaire et \textsf{booktabs}}
\label{sec:regles_booktabs}

Un tableau n'est pas une grille de tableur : c'est un dispositif typographique destiné à ordonner des données pour faciliter la comparaison intellectuelle.
Dans son manuel de référence devenu le standard mondial des revues à comité de lecture~\cite{fear2020booktabs}, Simon Fear énonce deux axiomes incontournables :

\begin{boiteSynthese}{Les deux axiomes de Simon Fear (paquet \textsf{booktabs})}
  \begin{enumerate}
    \item \textbf{Ne jamais utiliser de filets verticaux.} Les barres verticales entravent le balayage visuel horizontal de gauche à droite et créent un quadrillage lourd et inutile.
    \item \textbf{Ne jamais doubler les filets horizontaux.} Une structure claire s'obtient par la modulation de l'épaisseur des traits et l'équilibre des blancs, non par la multiplication de filets parasites.
  \end{enumerate}
\end{boiteSynthese}

\subsection{La hiérarchie des trois règles horizontales}
\label{subsec:trois_regles_horizontales}

Un tableau scientifique standard s'articule autour de trois filets de graisses distinctes :
\begin{itemize}
  \item \texttt{\textbackslash toprule} : Filet supérieur d'ouverture, d'une épaisseur marquée (généralement \qty{0.08}{\centi\meter}).
  \item \texttt{\textbackslash midrule} : Filet intermédiaire plus fin (\qty{0.05}{\centi\meter}), séparant l'en-tête explicatif des données numériques ou textuelles.
  \item \texttt{\textbackslash bottomrule} : Filet de fermeture d'épaisseur identique à \texttt{\textbackslash toprule}.
  \item \texttt{\textbackslash cmidrule(lr)\{i-j\}} : Filet sous-colonne partiel reliant les colonnes $i$ à $j$. L'option \texttt{(lr)} raccourcit subtilement les extrémités gauche et droite pour séparer nettement deux sous-groupes adjacents.
  \item \texttt{\textbackslash addlinespace} : Insère un espace vertical délicat entre deux blocs logiques de données, évitant de surcharger le tableau d'un trait continu.
\end{itemize}

\subsection{Positionnement normé des légendes}
\label{subsec:position_legendes_tableaux}

Contrairement aux figures dont la légende est systématiquement placée \textbf{en dessous} du graphique, la norme internationale stipule que \textbf{la légende d'un tableau doit impérativement être placée au-dessus} du tableau.
En effet, le lecteur doit prendre connaissance du contexte et de la signification des variables avant d'analyser les valeurs de la matrice.

\texttt{ClasseGuillaume} garantit cette convention grâce à l'option KOMA-Script \texttt{captions=tableheading} et au paramétrage du paquet \textsf{caption} :
\begin{verbatim}
\captionsetup[table]{position = above, skip = 8pt}
\end{verbatim}

\subsection{Démonstration comparative : Tableau de synthèse de mesures}
\label{subsec:demo_synthese_mesures}

Le tableau modulaire ci-dessous illustre l'usage conjugué de \textsf{booktabs}, \textsf{tabularx} et \textsf{siunitx} :

%------------------------------------------------------------------------------------
%  TABLEAU 1 : SYNTHÈSE DE MESURES EXPÉRIMENTALES (booktabs + tabularx + siunitx)
%------------------------------------------------------------------------------------
\begin{table}[htbp]
  \centering
  \small
  \caption{Synthèse des mesures thermomécaniques et validation expérimentale}
  \label{tab:synthese_mesures}
  \begin{tabularx}{\textwidth}{l c c c S[table-format=3.2] X}
    \toprule
    \textbf{Échantillon} & \textbf{Température} & \textbf{Pression} & \textbf{Vitesse} & {\textbf{Rendement (\%)}} & \textbf{Observations qualitatives} \\
    \midrule
    ECH-01 & \qty{22.4}{\degreeCelsius} & \qty{101.3}{\kilo\pascal} & \qty{1.45}{\meter\per\second} & 88.50 & Écoulement laminaire stable sans turbulence. \\
    ECH-02 & \qty{45.0}{\degreeCelsius} & \qty{150.2}{\kilo\pascal} & \qty{2.80}{\meter\per\second} & 92.15 & Amorçage de vortex en sortie de tuyère. \\
    ECH-03 & \qty{75.8}{\degreeCelsius} & \qty{210.0}{\kilo\pascal} & \qty{4.12}{\meter\per\second} & 95.40 & Régime turbulent pleinement développé. \\
    \addlinespace
    ECH-04 & \qty{90.2}{\degreeCelsius} & \qty{280.5}{\kilo\pascal} & \qty{5.60}{\meter\per\second} & 94.80 & Perte de charge par friction pariétale. \\
    ECH-05 & \qty{110.0}{\degreeCelsius} & \qty{350.0}{\kilo\pascal} & \qty{7.25}{\meter\per\second} & 91.20 & cavitation détectée au col d'étranglement. \\
    \bottomrule
  \end{tabularx}
\end{table}



Comme le met en évidence le~\cref{tab:synthese_mesures}, l'absence de traits verticaux allège considérablement la lecture et confère une élégance immédiate à la présentation des résultats.


%------------------------------------------------------------------------------------
%  SECTION 4.2 : TABLEAUX AVANCÉS EN SITUATION RÉELLE
%------------------------------------------------------------------------------------
\section{Tableaux Complexes, Multipage et Orientation Paysage}
\label{sec:tableaux_avances}

La mise en forme de données expérimentales denses exige des techniques avancées pour s'adapter à la largeur de page sans tronquer le texte et sans générer d'erreurs de dépassement (\texttt{overfull \textbackslash hbox}).

\subsection{Dimensionnement automatique avec \textsf{tabularx}}
\label{subsec:tabularx_auto}

En \LaTeX{} standard, l'environnement \texttt{tabular} requiert de fixer manuellement la largeur de chaque colonne (ex. \texttt{p\{4cm\}}).
Si le texte est trop long, il déborde de la marge de la page.
À l'inverse, l'environnement \texttt{tabularx} introduit le spécificateur de colonne majuscule \textbf{\texttt{X}} :
\begin{itemize}
  \item Il calcule automatiquement l'espace résiduel sur la largeur totale prescrite (\texttt{\textbackslash textwidth}).
  \item Il répartit proportionnellement cet espace entre toutes les colonnes \texttt{X}.
  \item Il justifie automatiquement le texte à l'intérieur de la cellule avec césure propre.
\end{itemize}

Le~\cref{tab:comparatif_classe} met en pratique cette flexibilité pour confronter les atouts de \texttt{ClasseGuillaume} face à une installation standard :

%------------------------------------------------------------------------------------
%  TABLEAU COMPARATIF : LATEX STANDARD VS CLASSEGUILLAUME
%------------------------------------------------------------------------------------
\begin{table}[htbp]
  \centering
  \small
  \caption{Comparatif synthétique des fonctionnalités : \LaTeX{} Standard vs \texttt{ClasseGuillaume}}
  \label{tab:comparatif_classe}
  \begin{tabularx}{\textwidth}{p{3.2cm} p{5.5cm} X}
    \toprule
    \textbf{Fonctionnalité} & \textbf{En \LaTeX{} Standard (\emph{Vanilla})} & \textbf{Avec \texttt{ClasseGuillaume.cls}} \\
    \midrule
    \textbf{Gestion des polices} & Déclarations manuelles \textsf{fontspec} avec fichiers physiques \texttt{.ttf} & Détection intelligente triple : système direct, \TeX{} Live natif ou repli Overleaf. \\
    \addlinespace
    \textbf{Interligne des titres} & \texttt{\textbackslash onehalfspacing} étire affreusement les titres multi-lignes & Protection automatique : interligne 1,5 dans le corps, interligne simple strict dans les titres. \\
    \addlinespace
    \textbf{Puces de paragraphes} & \texttt{\textbackslash paragraph} renvoie à la ligne ou nécessite un réglage KOMA & Puce colorée en ligne prête à l'emploi (\texttt{\textbackslash textbullet~}) avec poursuite immédiate du texte. \\
    \addlinespace
    \textbf{Changement de marges} & Écriture manuelle lourde de \texttt{\textbackslash newgeometry} à chaque bascule & Trois commandes immédiates : \texttt{\textbackslash geometrieWord}, \texttt{\textbackslash geometrieLarge}, \texttt{\textbackslash geometrieStandard}. \\
    \addlinespace
    \textbf{Notes de révision} & \textsf{todonotes} déborde des marges, casse en UTF-8 et sans infobulle & \texttt{\textbackslash todoH} : surlignage \textsf{lua-ul}, troncature automatique > 30 car., infobulle PDF et lien retour. \\
    \addlinespace
    \textbf{Illustrations dédiées} & Schémas et photos se mélangent aux figures dans la table des figures & Flottant \texttt{illustration} avec sa propre liste indépendante (\texttt{\textbackslash listofillustrations}). \\
    \addlinespace
    \textbf{Mise en valeur visuelle} & Environnements bruts ou code \textsf{tcolorbox} complexe à réécrire & Boîtes \texttt{cadreFigure} (coins arrondis) et \texttt{boiteSynthese} (fond teinté) prêtes à l'emploi. \\
    \addlinespace
    \textbf{Logo institutionnel} & Montage artisanal de paquets d'en-tête sujet aux décalages & Couche dynamique \textsf{scrlayer} avec bascules \texttt{\textbackslash masquerLogo} et \texttt{\textbackslash afficherLogo}. \\
    \addlinespace
    \textbf{Bibliographie FR} & Chaînes en anglais (\emph{visited on}, \emph{available from}) & Localisation soignée : « consulté le », « disponible sur », liens de retour \texttt{backref}. \\
    \bottomrule
  \end{tabularx}
\end{table}



\subsection{Gestion des grands tableaux multipages avec \textsf{longtable}}
\label{subsec:longtable_multipage}

Lorsqu'un inventaire de composants ou une série temporelle de mesures s'étend sur plusieurs pages consécutives, l'environnement flottant \texttt{table} ne peut pas se scinder.
Il faut faire appel au paquet \textsf{longtable} :

\begin{boiteSynthese}{Les quatre zones clés d'un tableau \textsf{longtable}}
  Un tableau multipage professionnel doit répéter ses en-têtes pour que le lecteur ne perde jamais le fil :
  \begin{itemize}
    \item \texttt{\textbackslash endfirsthead} : En-tête principal affiché sur la toute première page.
    \item \texttt{\textbackslash endhead} : En-tête secondaire répété en haut de chaque page suivante avec la mention explicite « \textit{Suite de la page précédente} ».
    \item \texttt{\textbackslash endfoot} : Pied de tableau affiché au bas des pages intermédiaires (« \textit{Suite sur la page suivante...} »).
    \item \texttt{\textbackslash endlastfoot} : Pied de clôture définitif arborant la règle \texttt{\textbackslash bottomrule}.
  \end{itemize}
\end{boiteSynthese}

Le~\cref{tab:composants_multipage} démontre l'intégration complète de cette structure :

%------------------------------------------------------------------------------------
%  TABLEAU MULTIPAGE AVEC LONGTABLE ET BOOKTABS
%------------------------------------------------------------------------------------
\begin{small}
\begin{longtable}{p{2.5cm} p{3.2cm} p{4.5cm} p{3.5cm}}
  \caption{Inventaire exhaustif des composants expérimentaux (Tableau multipage)}
  \label{tab:composants_multipage} \\
  \toprule
  \textbf{Référence} & \textbf{Désignation} & \textbf{Spécifications Techniques} & \textbf{Statut Métrologique} \\
  \midrule
  \endfirsthead

  \multicolumn{4}{c}{\tablename\ \thetable\ -- \textit{Suite de la page précédente}} \\
  \toprule
  \textbf{Référence} & \textbf{Désignation} & \textbf{Spécifications Techniques} & \textbf{Statut Métrologique} \\
  \midrule
  \endhead

  \midrule
  \multicolumn{4}{r}{\textit{Suite sur la page suivante \dots}} \\
  \bottomrule
  \endfoot

  \bottomrule
  \endlastfoot

  CMP-101 & Capteur piézoélectrique & Plage \qtyrange{0}{50}{\kilo\pascal}, précision \qty{0.1}{\percent} & Étalonné ISO 17025 \\
  CMP-102 & Sonde thermocouple K & Plage \qtyrange{-50}{400}{\degreeCelsius}, temps de réponse \qty{20}{\milli\second} & Certifié d'usine \\
  CMP-103 & Débitmètre vortex & Diamètre intérieur \qty{25}{\milli\meter}, sortie \qtyrange{4}{20}{\milli\ampere} & Vérification annuelle \\
  CMP-104 & Manomètre différentiel & Échelle \qtyrange{0}{5}{\bar}, membrane inox 316L & En attente d'ajustage \\
  \addlinespace
  CMP-201 & Pompe volumétrique & Cylindrée \qty{12.5}{\centi\meter\cubed}, moteur brushless & Validé banc d'essai \\
  CMP-202 & Vanne proportionnelle & Régulation PID intégrée, alimentation \qty{24}{\volt} DC & Zéro dérive mesuré \\
  CMP-203 & Échangeur à plaques & Surface d'échange \qty{1.8}{\meter\squared}, flux contre-courant & Test hydrostatique OK \\
  CMP-204 & Séparateur cyclonique & Efficacité \qty{99.2}{\percent} sur particules $> \qty{5}{\micro\meter}$ & Conforme norme ATEX \\
  \addlinespace
  CMP-301 & Module d'acquisition & 16 canaux analogiques, conversion 24 bits à \qty{100}{\kilo\hertz} & Horloge synchronisée \\
  CMP-302 & Convertisseur USB-CAN & Isolation galvanique \qty{2.5}{\kilo\volt}, débit \qty{1}{\mega\bit\per\second} & Pilote validé Linux \\
  CMP-303 & Alimentation stabilisée & Tension \qtyrange{0}{30}{\volt}, courant maxi \qty{10}{\ampere} & Ondulation $< \qty{1}{\milli\volt} $ \\
  CMP-304 & Blindage électromagnétique & Atténuation \qty{65}{\decibel} entre \qty{10}{\mega\hertz} et \qty{1}{\giga\hertz} & Cage de Faraday testée \\
\end{longtable}
\end{small}



\subsection{Tableaux panoramiques en mode paysage avec \textsf{pdflscape}}
\label{subsec:mode_paysage}

Certaines synthèses transversales comportent un nombre trop élevé de paramètres pour être contenues dans la largeur portrait d'une feuille A4 (\qty{21}{\centi\meter}).
L'utilisation de l'environnement \texttt{landscape} issu du paquet \textsf{pdflscape} résout cette contrainte avec un confort de lecture total :
contrairement au paquet classique \textsf{lscape} qui laisse la page en portrait tout en tournant le texte de 90 degrés, \textsf{pdflscape} transmet une directive au visualiseur PDF pour effectuer une \textbf{véritable rotation physique de la page}.
Le document reste ainsi parfaitement lisible à l'écran sans que l'utilisateur n'ait à pivoter son moniteur.

Le~\cref{tab:evaluation_paysage} (page suivante) offre un exemple probant de matrice décisionnelle en mode paysage.

%------------------------------------------------------------------------------------
%  TABLEAU PAYSAGE LARGE (pdflscape + booktabs + tabularx)
%------------------------------------------------------------------------------------
\begin{landscape}
\begin{table}[p]
  \centering
  \small
  \caption{Matrice d'évaluation comparative multicritères des architectures logicielles et typographiques}
  \label{tab:evaluation_paysage}
  \begin{tabularx}{\linewidth}{p{3.5cm} p{3.5cm} c c c c X}
    \toprule
    \textbf{Moteur / Chaîne} & \textbf{Standard de Polices} & \textbf{UTF-8 Natif} & \textbf{Microtypo} & \textbf{Vitesse} & \textbf{Complexité} & \textbf{Conclusion Opérationnelle} \\
    \midrule
    pdf\LaTeX{} classique & T1 / Type 1 (8 bits) & Non (inputenc) & Limitée & Rapide & Élevée (bidouilles) & Obsolète pour nouveaux projets académiques. \\
    \addlinespace
    Xe\LaTeX{} & OpenType / TrueType & Oui & Basique & Moyenne & Moyenne & Préféré historiquement pour langues non latines. \\
    \addlinespace
    \textbf{Lua\LaTeX{} moderne} & \textbf{OpenType / TrueType} & \textbf{Oui (100\,\%)} & \textbf{Complète} & \textbf{Excellente} & \textbf{Épurée} & \textbf{Standard contemporain recommandé par l'UCLouvain.} \\
    \midrule
    \multicolumn{7}{l}{\textbf{Bilan comparatif des fonctionnalités de publication :}} \\
    \addlinespace
    Classe KOMA \texttt{scrbook} & Système KOMA-Script & Oui & Oui & -- & Modérée & Socle de base solide mais sans encadrements prêts à l'emploi. \\
    \addlinespace
    \texttt{ClasseGuillaume} & \textbf{Architecture Modulaire} & \textbf{Oui} & \textbf{Pénalités max} & -- & \textbf{Immédiate} & \textbf{Intègre \texttt{\textbackslash todoH}, 3 géométries, logo scrlayer et booktabs.} \\
    \bottomrule
  \end{tabularx}
  \vspace{4mm}
  \begin{minipage}{\linewidth}
    \footnotesize
    \textit{Note méthodologique :} Ce tableau panoramique démontre l'utilisation conjointe de l'environnement \texttt{landscape} fourni par le paquet \textsf{pdflscape} et de la classe de document. La rotation est physique dans le fichier PDF généré, garantissant une consultation parfaite sans devoir tourner manuellement l'écran.
  \end{minipage}
\end{table}
\end{landscape}





%====================================================================================
%  PARTIE III : DYNAMIQUE DE PAGE, MISE EN PAGE COMPLEXE & GRAPHISME
%====================================================================================
\part{Dynamique de Page, Mise en Page Complexe \& Graphisme}
\label{part:graphisme_geometries}

\chapter{Figures, Illustrations et Boîtes Décoratives}
\label{chap:illustrations_boites}
%------------------------------------------------------------------------------------
%  SECTION 5.1 : FLOTTANTS CLASSIQUES VS TYPE D'ILLUSTRATION DÉDIÉ
%------------------------------------------------------------------------------------
\section{Flottants Standards et Environnement Dédié aux Illustrations}
\label{sec:figures_illustrations}

La restitution graphique d'expériences ou d'analyses conceptuelles s'articule autour de deux types de supports visuels :
les figures techniques (graphiques vectoriels, courbes, schémas de principe) d'une part, et les illustrations iconographiques (clichés photographiques, reproductions anatomiques, œuvres artistiques) d'autre part.
Mélanger ces deux catégories au sein d'une unique « table des figures » nuit à la clarté du document.

\subsection{Figure standard avec sous-figures coordonnées (\textsf{subcaption})}
\label{subsec:sous_figures}

Pour juxtaposer deux résultats comparatifs, le paquet moderne \textsf{subcaption} remplace avantageusement le paquet obsolète \texttt{subfig}.
Chaque sous-figure est isolée dans un environnement \texttt{subfigure} avec sa propre sous-légende indicée \texttt{(a)}, \texttt{(b)} :

\begin{figure}[htbp]
  \centering
  \begin{subfigure}[b]{0.46\textwidth}
    \centering
    \begin{tcolorbox}[colframe=secondaryColor, colback=white, boxrule=1pt, arc=4pt, width=\linewidth, halign=center]
      \vspace{0.8cm}
      {\large\sffamily\bfseries\color{secondaryColor} [ Schéma Technique A ]}\\[2mm]
      {\footnotesize Flux d'injection sous pression \qty{150}{\kilo\pascal}}
      \vspace{0.8cm}
    \end{tcolorbox}
    \caption{Régime d'injection stationnaire}
    \label{subfig:injection_stationnaire}
  \end{subfigure}
  \hfill
  \begin{subfigure}[b]{0.46\textwidth}
    \centering
    \begin{tcolorbox}[colframe=secondaryColor, colback=white, boxrule=1pt, arc=4pt, width=\linewidth, halign=center]
      \vspace{0.8cm}
      {\large\sffamily\bfseries\color{secondaryColor} [ Schéma Technique B ]}\\[2mm]
      {\footnotesize Propagation d'ondes à \qty{4.2}{\meter\per\second}}
      \vspace{0.8cm}
    \end{tcolorbox}
    \caption{Régime transitoire perturbé}
    \label{subfig:injection_transitoire}
  \end{subfigure}
  \caption{Comparaison des dynamiques d'écoulement sous deux régimes de fonctionnement}
  \label{fig:comparaison_regimes}
\end{figure}

Comme le démontre la~\cref{fig:comparaison_regimes}, et plus spécifiquement la~\cref{subfig:injection_stationnaire}, chaque sous-figure peut être référencée indépendamment dans le corps du texte grâce à \textsf{cleveref}.

\subsection{L'innovation de \texttt{ClasseGuillaume} : Le flottant \texttt{illustration}}
\label{subsec:flottant_illustration}

Pour répondre aux standards des grandes monographies et thèses universitaires, la classe déclare un nouveau type de flottant officiel via le moteur \textsf{tocbasic} de KOMA-Script :

\begin{verbatim}
\DeclareNewTOC[%
  type = illustration,
  float,
  name = Illustration,
  listname = {Liste des illustrations},
]{loi}
\end{verbatim}

Cet environnement s'utilise exactement comme une figure ordinaire, mais bénéficie d'une numérotation propre et alimente une \textbf{table des illustrations indépendante} (\texttt{\textbackslash listofillustrations}) située en fin d'ouvrage :

\begin{illustration}[htbp]
  \centering
  \begin{cadreFigure}[width=0.82\textwidth]
    \vspace{0.6cm}
    {\Large\sffamily\bfseries\color{primaryColor} Cliché Iconographique de Microscopie Électronique}\\[3mm]
    {\normalsize Grossissement $\times \num{25000}$ -- Structure cristalline des nanotubes de carbone}
    \vspace{0.6cm}
  \end{cadreFigure}
  \caption{Observation nanoscopique d'échantillons en champ sombre}
  \label{ill:microscopie_nanotubes}
\end{illustration}

L'\cref{ill:microscopie_nanotubes} est reconnue de manière transparente par \textsf{cleveref}, qui lui attribue automatiquement le préfixe « illustration » en minuscules ou « Illustration » en majuscules en début de phrase.
L'ensemble de ces supports est répertorié dans la \emph{Liste des illustrations} du Post-Contenu.


%------------------------------------------------------------------------------------
%  SECTION 5.2 : VALORISATION VISUELLE AVEC LES BOÎTES TCOLORBOX ET TIKZ
%------------------------------------------------------------------------------------
\section{Encadrés Typographiques Modernes et Schémas Vectoriels TikZ}
\label{sec:boites_tcolorbox_tikz}

L'aération d'un document passe par l'insertion rythmée d'éléments d'accroche visuelle.
En \LaTeX{} classique, la création d'encadrés personnalisés imposait des constructions complexes à base de \texttt{\textbackslash fbox} ou \texttt{mdframed}.
La classe \texttt{ClasseGuillaume} intègre le moteur de référence \textsf{tcolorbox} (avec les extensions \texttt{skins} et \texttt{breakable}), offrant des boîtes insensibles aux sauts de page et dotées d'une finition graphique irréprochable.

\subsection{L'environnement \texttt{cadreFigure} pour sublimer les schémas}
\label{subsec:cadre_figure}

L'environnement \texttt{cadreFigure} habille son contenu d'un contour raffiné aux coins délicatement arrondis (\texttt{arc = 10pt}) et d'un liseré de couleur primaire assombrie :

\begin{figure}[htbp]
  \centering
  \begin{cadreFigure}[width=0.92\textwidth]
    \begin{tikzpicture}[
      node distance = 2.1cm,
      every node/.style = {font=\small\sffamily},
      block/.style = {draw=primaryColor, thick, fill=primaryColor!8, rounded corners=3pt, minimum width=2.4cm, minimum height=1.1cm, align=center},
      sum/.style = {draw=secondaryColor, thick, circle, fill=secondaryColor!10, minimum size=0.8cm},
      arrow/.style = {->, >=stealth, thick, color=darkgray}
    ]
      % Noeuds du schéma bloc
      \node[sum] (sommateur) {$\Sigma$};
      \node[block, right of=sommateur, xshift=1.2cm] (correcteur) {Correcteur\\Numérique $C(z)$};
      \node[block, right of=correcteur, xshift=1.6cm] (processus) {Procédé Physique\\Modélisé $G(s)$};
      \node[coordinate, right of=processus, xshift=1.2cm] (sortie) {};
      \node[coordinate, below of=correcteur, yshift=0.8cm] (retour) {};

      % Flèches de connexion
      \draw[arrow] ([xshift=-1.6cm]sommateur.west) -- node[above] {$r(t)$} node[below, pos=0.2] {\footnotesize Consigne} (sommateur.west);
      \draw[arrow] (sommateur) -- node[above] {$e(t)$} (correcteur);
      \draw[arrow] (correcteur) -- node[above] {$u(t)$} (processus);
      \draw[arrow] (processus) -- node[above] {$y(t)$} (sortie);
      \draw[arrow] ([xshift=0.8cm]processus.east) |- (retour) -| node[pos=0.9, left] {$-$} (sommateur.south);
    \end{tikzpicture}
  \end{cadreFigure}
  \caption{Architecture fonctionnelle d'une boucle d'asservissement en temps réel composée sous TikZ}
  \label{fig:schema_bloc_tikz}
\end{figure}

Le schéma vectoriel représenté à la~\cref{fig:schema_bloc_tikz} illustre la symbiose parfaite entre le graphisme vectoriel natif \textsf{TikZ} et le conteneur \texttt{cadreFigure}.

\subsection{L'environnement \texttt{boiteSynthese} pour les messages clés}
\label{subsec:boite_synthese}

L'environnement \texttt{boiteSynthese\{Titre\}} structure les conclusions d'étapes, définitions ou théorèmes au moyen d'un fond teinté doux (\texttt{primaryColor!6}) et d'un titre en lettres sans-serif grasses :

\begin{boiteSynthese}{Principe Fondamental de l'Équilibre Graphique}
  Les boîtes décoratives doivent soutenir la transmission intellectuelle et non la distraire :
  \begin{itemize}
    \item Une boîte ne doit contenir qu'une seule idée forte ou synthèse transversale.
    \item L'intensité de la couleur d'arrière-plan ne doit jamais dépasser \qty{8}{\percent} pour garantir un ratio de contraste d'impression optimal et éviter la fatigue oculaire.
    \item Dans un mémoire volumineux, limitez l'usage à une boîte de synthèse par section majeure.
  \end{itemize}
\end{boiteSynthese}

\subsection{Boîtes d'alerte avec bandeau vertical (\emph{Left-bar alert})}
\label{subsec:boite_alerte}

Pour les mises en garde méthodologiques ou les avertissements de sécurité expérimentale, un style minimaliste avec filet latéral (\emph{leftbar}) s'avère particulièrement sobre et percutant :

\begin{tcolorbox}[
  enhanced,
  blanker,
  borderline west={3pt}{0pt}{Bittersweet},
  left=12pt,
  top=4pt,
  bottom=4pt
]
  \textcolor{Bittersweet}{\bfseries\sffamily Règle de Sécurité et Hygiène des Données :}\\
  \small\itshape Lors de l'acquisition de signaux physiologiques sur sujets humains, assurez-vous impérativement de la conformité avec le Règlement Général sur la Protection des Données (RGPD) et de l'anonymisation irréversible des identifiants patients avant tout traitement statistique.
\end{tcolorbox}



\chapter{Géométries Multiples, Styles de Pages, Colonnes \& Logo Dynamique}
\label{chap:geometries_styles}
%------------------------------------------------------------------------------------
%  SECTION 6.1 : GÉOMÉTRIES MULTIPLES À LA VOLÉE
%------------------------------------------------------------------------------------
\section{Alternance Dynamique des Géométries de Page}
\label{sec:geometries_multiples}

Dans un mémoire ou un rapport d'ingénierie, une taille de marge unique pour l'ensemble du volume est inadaptée :
la page de couverture et les grands tableaux requièrent des marges étroites pour maximiser l'espace utile, le corps du texte nécessite une largeur de colonne ergonomique (65 à 75 caractères par ligne pour un confort visuel optimal), tandis que la phase de relecture collaborative exige une marge latérale généreuse pour accueillir les annotations.

En \LaTeX{} standard, changer de géométrie en cours de document provoque souvent des décalages imprévisibles dans les en-têtes ou rompt l'alignement recto-verso.
\texttt{ClasseGuillaume} résout cette difficulté en encapsulant la gestion du paquet \textsf{geometry} et du moteur \textsf{typearea} de KOMA-Script au sein de trois commandes directes et sécurisées.

\subsection{Les trois commandes de géométrie de la classe}
\label{subsec:trois_commandes_geometrie}

\begin{itemize}
  \item \textbf{\texttt{\textbackslash geometrieWord} :}
  Fixe des marges symétriques de \qty{2.5}{\centi\meter} sur les quatre bords et alloue une marge latérale dédiée de \qty{2.2}{\centi\meter} (\texttt{marginparwidth}) avec un séparateur de \qty{3}{\milli\meter} (\texttt{marginparsep}).
  C'est la géométrie par excellence pour le corps de texte rédigé et la phase de relecture collaborative avec les notes \texttt{\textbackslash todoH}.
  
  \item \textbf{\texttt{\textbackslash geometrieLarge} (ou \texttt{\textbackslash geometrieExemple}) :}
  Applique des marges compactes de \qty{1.2}{\centi\meter} à gauche et à droite, et \qty{1.6}{\centi\meter} en haut et en bas.
  Elle annule l'espace marginal pour offrir une largeur utile maximale de \qty{18.6}{\centi\meter}, idéale pour la page de garde, les posters de synthèse, ou les grands tableaux récapitulatifs.

  \item \textbf{\texttt{\textbackslash geometrieStandard} :}
  Restaure instantanément la mise en page canonique calculée par l'algorithme typographique de KOMA-Script (\texttt{DIV = 12}), garantissant des proportions harmonieuses selon les canons du livre d'art.
\end{itemize}

\begin{boiteSynthese}{Règle de bascule propre}
  Chacune de ces commandes intègre automatiquement une directive \texttt{\textbackslash cleardoublepage} avant d'appliquer les nouvelles marges.
  Cette précaution fondamentale garantit qu'une nouvelle géométrie commence toujours sur une page neuve et prévient tout débordement de texte entre deux pages consécutives aux proportions divergentes.
\end{boiteSynthese}


%------------------------------------------------------------------------------------
%  SECTION 6.2 : EN-TÊTES ET PIEDS DE PAGE MODULAIRES AVEC SCRLAYER-SCRPAGE
%------------------------------------------------------------------------------------
\section{Styles d'En-têtes et Pieds de Page Modulaires}
\label{sec:styles_en_tetes_pieds}

La gestion des hauts et bas de page en \LaTeX{} a longtemps été confiée au paquet obsolète \texttt{fancyhdr}.
Dans un document moderne basé sur KOMA-Script, \texttt{fancyhdr} provoque des conflits majeurs avec la géométrie et le placement des notes de marge.
\texttt{ClasseGuillaume} s'appuie sur le système de calques haute fidélité \textsf{scrlayer-scrpage}, offrant une modularité sans précédent.

\subsection{Le style par défaut : \texttt{scrheadings}}
\label{subsec:style_scrheadings}

Le style actif par défaut dans l'ensemble du mémoire est \texttt{scrheadings} :
\begin{itemize}
  \item Un filet horizontal fin (\qty{0.5}{pt}) sépare l'en-tête du texte courant.
  \item L'en-tête intérieur affiche dynamiquement le titre du chapitre ou de la section courante via la macro \texttt{\textbackslash headmark}.
  \item L'en-tête extérieur affiche le numéro de page en chiffres gras (\texttt{\textbackslash pagemark}).
  \item Le pied de page reste volontairement épuré pour ne pas alourdir la composition.
\end{itemize}

\subsection{Démonstration en direct du style épuré : \texttt{StyleFocus}}
\label{subsec:style_focus}

Pour les lectures concentrées, les synthèses de fin de partie ou les rapports managériaux, la classe fournit le style \texttt{StyleFocus}.
Ce style supprime le filet d'en-tête, centre le rappel de section entre tirets (\texttt{-- \textbackslash headmark{} --}) et positionne le numéro de page au centre du pied de page.

\begingroup
% Bascule locale en StyleFocus
\pagestyle{StyleFocus}

Nous activons immédiatement \texttt{StyleFocus} sur cette portion de texte.
Observez l'en-tête et le pied de cette page :
la ligne supérieure a disparu, laissant place à une respiration minimale, tandis que la pagination est descendue au centre du pied de page.
Cette disposition est particulièrement prisée pour les sections de réflexion ou les déclarations éthiques.

\subsection{Démonstration en direct du style de diagnostic : \texttt{StyleTest}}
\label{subsec:style_test}

Durant la phase de développement et de révision technique, l'auteur a besoin d'informations de débogage précises : quel est le nom du fichier compilé ? Combien de pages totales compte le document ?

La classe intègre à cet effet le style technique \texttt{StyleTest} :
\begin{itemize}
  \item Un badge orange centré \texttt{\tiny [DEBUG VIEW]} dans l'en-tête.
  \item Des filets supérieur (\qty{0.8}{pt}) et inférieur (\qty{0.4}{pt}).
  \item L'indication explicite du fichier racine (\texttt{\textbackslash jobname.tex}) dans le pied de page.
  \item Le décompte complet des pages via le paquet \textsf{lastpage} : \texttt{p.~\textbackslash thepage/\textbackslash pageref\{LastPage\}}.
\end{itemize}

% Bascule locale en StyleTest
\pagestyle{StyleTest}

Cette nouvelle bascule active instantanément la vue de diagnostic sur les pages courantes.
Elle permet aux équipes de relecture de vérifier en un coup d'œil l'horodatage et la pagination relative sans devoir inspecter les métadonnées internes du fichier PDF.

\endgroup
% Rétablissement automatique de scrheadings à la sortie du groupe
\pagestyle{scrheadings}

La sortie du bloc rétablit automatiquement le style standard \texttt{scrheadings} pour la suite du mémoire.
Ces bascules illustrent la flexibilité extrême offerte par la classe pour rythmer la lecture selon la nature des contenus.


%------------------------------------------------------------------------------------
%  SECTION 6.3 : MISES EN PAGE MULTI-COLONNES (2 ET 3 COLONNES)
%------------------------------------------------------------------------------------
\section{Mises en Page Multi-Colonnes et Équilibrage Dynamique}
\label{sec:multi_colonnes}

Pour les documents techniques denses, les synthèses de laboratoire ou les articles de revue, une disposition en colonne unique (\emph{single-column}) peut engendrer des lignes trop longues (plus de 85 caractères), rendant la lecture fatigante.
Le passage en colonnes multiples permet de densifier l'information tout en conservant une longueur de ligne idéale (40 à 50 caractères).

\subsection{L'option globale \texttt{twocolumn} vs la flexibilité locale de \textsf{multicol}}
\label{subsec:comparatif_colonnes}

Il existe deux manières d'aborder le multi-colonnes en \LaTeX{} :
\begin{enumerate}
  \item \textbf{L'option globale de classe (\texttt{\textbackslash documentclass[twocolumn]\{ClasseGuillaume\}}) :}
  Elle scinde l'ensemble du mémoire en deux colonnes du début à la fin.
  Ce mode convient aux formats d'articles IEEE ou de comptes-rendus de conférences, mais s'avère trop rigide pour un mémoire universitaire dont la page de garde, les résumés et les grands tableaux doivent rester en pleine largeur.
  
  \item \textbf{L'approche locale dynamique avec le paquet \textsf{multicol} :}
  Elle permet de basculer ponctuellement en deux, trois ou quatre colonnes sur une fraction de page, puis de revenir à la colonne pleine de façon parfaitement fluide et sans saut de page forcé.
\end{enumerate}

\subsection{Démonstration en direct d'un passage en deux colonnes équilibrées}
\label{subsec:demo_deux_colonnes}

Voici une démonstration concrète d'un texte scindé en deux colonnes avec filet vertical séparateur délicat :

\setlength{\columnsep}{22pt}
\setlength{\columnseprule}{0.4pt}

\begin{multicols}{2}[\subsubsection*{Étude comparative des fluides caloporteurs en régime forcé}]
  Dans cette première colonne, nous examinons les propriétés thermophysiques de l'eau glycolée à \qty{30}{\percent}.
  Sa capacité thermique massique élevée ($c_p = \qty{3.85}{\kilo\joule\per\kilo\gram\per\kelvin}$) en fait un fluide de choix pour les circuits de refroidissement industriels soumis à des températures négatives modérées.
  
  Toutefois, sa viscosité dynamique augmente de manière sensible aux basses températures, entraînant une majoration des pertes de charge singulières dans les coudes de canalisations.
  
  \columnbreak % Forçage pédagogique de l'équilibrage de colonne
  
  Dans cette seconde colonne en regard, nous analysons le comportement d'une huile synthétique de synthèse.
  Bien que sa conductivité thermique soit inférieure ($\lambda = \qty{0.13}{\watt\per\meter\per\kelvin}$), sa très haute stabilité thermique jusqu'à \qty{320}{\degreeCelsius} autorise un fonctionnement sans pressurisation excessive du circuit primaire.
  
  L'équilibrage automatique assuré par le paquet \textsf{multicol} égalise précisément la hauteur des deux colonnes avant de restituer la largeur normale de page.
\end{multicols}

\setlength{\columnseprule}{0pt} % Rétablissement du séparateur invisible

\subsection{Démonstration d'un lexique technique en trois colonnes}
\label{subsec:demo_trois_colonnes}

Pour les nomenclatures, index de variables ou listes de constantes, le format en trois colonnes optimise l'occupation de la page :

\setlength{\columnsep}{16pt}
\begin{multicols}{3}
  \begin{itemize}[leftmargin=12pt]
    \item $\alpha$ : Diffusivité thermique
    \item $\beta$ : Dilatation volumique
    \item $\gamma$ : Indice adiabatique
    \item $\delta$ : Épaisseur de couche
    \item $\epsilon$ : Émissivité radiative
    \item $\zeta$ : Coefficient d'amortissement
    \item $\eta$ : Rendement isentropique
    \item $\theta$ : Température adimensionnelle
    \item $\kappa$ : Compressibilité isotherme
    \item $\lambda$ : Conductivité thermique
    \item $\mu$ : Viscosité dynamique
    \item $\nu$ : Viscosité cinématique
  \end{itemize}
\end{multicols}


%------------------------------------------------------------------------------------
%  SECTION 6.4 : LOGO DYNAMIQUE SCRLAYER ET DOCUMENTS PDF EXTERNES
%------------------------------------------------------------------------------------
\section{Logo en Calque Dynamique et Insertion de Documents Externes}
\label{sec:logo_documents_externes}

Deux défis récurrents dans la formalisation d'un livrable officiel concernent la présence d'un filigrane institutionnel discret et la fusion de pièces jointes administratives ou techniques fournies par des tiers.

\subsection{Le filigrane dynamique en arrière-plan (\textsf{scrlayer})}
\label{subsec:calque_logo}

L'insertion manuelle d'un logo dans les en-têtes (\texttt{\textbackslash ihead} ou \texttt{\textbackslash ohead}) modifie l'encombrement vertical (\texttt{headheight}) et crée des décalages de texte selon que la page est paire ou impaire.

\texttt{ClasseGuillaume} s'affranchit de cette contrainte en utilisant le concept de **couche d'arrière-plan** (\emph{background layer}) du moteur \textsf{scrlayer} :

\begin{verbatim}
\DeclareNewLayer[
  background,
  topmargin,
  align = tr,
  hoffset = \paperwidth - 1.8cm, % Calibré dynamiquement sur le bord droit
  voffset = 1.2cm,
  width = 3cm,
  height = 3cm,
  contents = {%
    \IfFileExists{images/logo_test.png}{%
      \includegraphics[width=2.2cm]{images/logo_test.png}%
    }{}%
  }
]{logo.layer}
\end{verbatim}

\begin{boiteSynthese}{Commandes de pilotage du logo institutionnel}
  Le calque flotte au-dessus de la mise en page sans jamais interférer avec le texte du corps. Deux commandes permettent d'en moduler l'affichage selon les chapitres :
  \begin{itemize}
    \item \texttt{\textbackslash masquerLogo} : Désactive instantanément le calque (utilisé dans le Frontmatter pour préserver la virginité de la page de titre et du résumé).
    \item \texttt{\textbackslash afficherLogo} : Réactive le calque en tête du corps de texte (\texttt{\textbackslash mainmatter}) ou dans les sections spécialisées.
  \end{itemize}
\end{boiteSynthese}

\subsection{Intégration transparente de documents PDF externes avec \textsf{pdfpages}}
\label{subsec:pdfpages_externes}

Dans une thèse ou un rapport d'ingénierie, il est fréquent de devoir intégrer des documents tiers au format PDF :
attestations administratives, comptes-rendus de laboratoires accrédités, fiches techniques de constructeurs, ou articles scientifiques complets déjà publiés.
Réinsérer ces documents sous forme d'images capturées dégrade fortement la résolution et perd la numérotation séquentielle des pages.

La classe intègre le paquet de référence \textsf{pdfpages} :
\begin{verbatim}
% Exemple d'inclusion d'un document multipage en annexe :
\includepdf[
  pages = {1-3},                  % Pages à importer
  scale = 0.92,                   % Échelle réservant la marge KOMA
  pagecommand = {\thispagestyle{scrheadings}} % En-têtes conservés
]{documents/rapport_mesures.pdf}
\end{verbatim}
Cette méthode vectorielle conserve la netteté parfaite des textes et graphiques d'origine tout en assurant une continuité sans faille de la pagination générale de votre mémoire.
Une démonstration concrète est proposée en Annexe~\ref{annexe:specimen_externe}.



%====================================================================================
%  PARTIE IV : COLLABORATION, RÉFÉRENCEMENT & DIFFUSION
%====================================================================================
\part{Collaboration, Référencement \& Diffusion}
\label{part:collaboration_diffusion}

\chapter{Outils Collaboratifs, Système d'Annotation \texttt{\textbackslash todoH} \& Révision}
\label{chap:collaboration_todos}
%------------------------------------------------------------------------------------
%  SECTION 7.1 : LE SYSTÈME EXCLUSIF D'ANNOTATIONS COLLABORATIVES TODOH
%------------------------------------------------------------------------------------
\section{Le Système Exclusif d'Annotations Collaboratives \texttt{\textbackslash todoH}}
\label{sec:systeme_todoh}

La phase de relecture entre étudiants, directeurs de thèse et pairs académiques est une étape cruciale.
Le paquet standard \textsf{todonotes} est souvent employé, mais il souffre de trois limitations majeures :
les messages longs débordent verticalement du bas de page, les caractères accentués UTF-8 génèrent des erreurs de compilation avec les paquets de soulignement traditionnels, et le texte surligné dans le corps n'offre aucun lien interactif vers le tableau de bord des tâches.

La classe \texttt{ClasseGuillaume} résout l'ensemble de ces verrous technologiques grâce à sa commande exclusive \texttt{\textbackslash todoH}.

\subsection{Les cinq innovations majeures de la macro \texttt{\textbackslash todoH}}
\label{subsec:cinq_innovations_todoh}

La syntaxe générale de la commande se décompose en quatre arguments :
\begin{center}
  \texttt{\textbackslash todoH\{options\_style\}\{couleur\}\{message\_marge\}\{texte\_surligné\}}
\end{center}

\begin{enumerate}
  \item \textbf{Surlignage robuste via \textsf{lua-ul} et \textsf{luacolor} :}
  Le texte cible est surligné dans la couleur du relecteur. Contrairement aux paquets obsolètes (\textsf{soul}), le moteur Lua gère les césures de mots, les accents français et les équations sans aucun bogue de compilation.
  
  \item \textbf{Troncature automatique des messages longs ($> 30$ caractères) :}
  Pour éviter qu'un commentaire volumineux ne vienne chevaucher la note suivante ou déborder de la marge, la classe mesure la longueur de la chaîne avec \textsf{xstring} et abrège la note à 30 caractères suivis de points de suspension.

  \item \textbf{Infobulle interactive (\textsf{pdfcomment}) :}
  Grâce à la commande \texttt{\textbackslash pdftooltip}, l'intégralité du message reste disponible : dans votre lecteur PDF (Adobe Acrobat, PDF-XChange, navigateurs modernes), \textbf{un simple survol de la souris sur la note dans la marge fait apparaître une bulle contextuelle} dévoilant le texte complet !

  \item \textbf{Détection intelligente des commandes et URL (\textsf{xurl}) :}
  Si le commentaire contient une balise \LaTeX{} (comme un lien hypertexte \texttt{\textbackslash url\{...\}}), la classe détecte la présence de la barre oblique inversée (\texttt{\textbackslash}) et suspend la troncature pour préserver l'intégrité du lien cliquable.

  \item \textbf{Lien hypertexte de retour vers la table des tâches :}
  Un clic direct sur le texte surligné dans le paragraphe téléporte instantanément le lecteur vers la liste récapitulative des révisions (\texttt{\textbackslash listoftodos}) située au début du mémoire.
\end{enumerate}

\subsection{Démonstrations concrètes en situation de révision}
\label{subsec:demos_todoh}

Voici deux exemples concrets illustrant le comportement adaptatif de la commande :

Dans ce premier passage, nous simulons une remarque de fond \todoH{FloraStyle}{red!30}{Cette démonstration théorique requiert une justification expérimentale plus détaillée avec des données chiffrées issues du banc d'essai.}{un raisonnement nécessitant un approfondissement méthodologique}.
Observez dans la marge de droite : la note est abrégée pour ne pas saturer l'espace, mais un survol à la souris révèle l'ensemble de la directive de relecture.

Dans ce second cas, nous transmettons à l'auteur une ressource bibliographique en ligne sans risque d'altération du lien : \todoH{GuillaumeStyle}{NavyBlue!30}{Vérifier les spécifications normatives sur \url{https://www.iso.org/standard/30669.html}}{la conformité avec les recommandations internationales ISO}.
Le lien URL est protégé et césuré proprement grâce au paquet \textsf{xurl}.


%------------------------------------------------------------------------------------
%  SECTION 7.2 : PROFILS DE RELECTEURS ET GESTION DES FIGURES PROVISOIRES
%------------------------------------------------------------------------------------
\section{Profils d'Auteurs Multiples et Emplacements Réservés}
\label{sec:profils_relecteurs}

La relecture d'un travail de recherche implique souvent plusieurs intervenants (promoteur, co-promoteur, expert externe, relecteur grammatical).
Chacun doit pouvoir être identifié immédiatement par sa couleur et son cartouche d'auteur.

\subsection{Les styles de relecteurs préconfigurés}
\label{subsec:styles_relecteurs_preconfig}

La classe intègre plusieurs styles de relecteurs prêts à l'emploi :
\begin{itemize}
  \item \texttt{GuillaumeStyle} : Attribué au rédacteur principal (\texttt{author=Guillaume}), agrémenté d'une ombre portée subtile.
  \item \texttt{MarieStyle} : Style dédié pour les remarques de direction académique (\texttt{author=Marie}).
  \item \texttt{FloraStyle} : Style dédié pour les corrections de forme ou méthodologiques (\texttt{author=Flora}).
  \item \texttt{TestStyle} : Style neutre pour les auto-vérifications techniques.
\end{itemize}

Démonstration croisée : ici, Marie note que \todoH{MarieStyle}{Emerald!30}{Penser à insérer le protocole de calibration des capteurs avant soutenance finale.}{cette étape de mesure devra être corrélée} avec les résultats théoriques obtenus au Chapitre~\ref{chap:sciences_maths}.

\subsection{Réservation de figures provisoires avec \texttt{\textbackslash missingfigure}}
\label{subsec:missing_figures_demo}

Durant la rédaction active, il est fréquent que certains graphiques expérimentaux soient encore en cours de calcul ou que des photographies de laboratoire ne soient pas finalisées.
Stopper la rédaction en attendant le fichier image briserait l'élan de travail.

La commande \texttt{\textbackslash missingfigure} réserve l'espace dans la mise en page sans bloquer la chaîne de compilation :

\begin{figure}[htbp]
  \centering
  \missingfigure[
    figwidth = 0.88\textwidth,
    figheight = 3.5cm,
    figcolor = RoyalBlue!15
  ]{Schéma vectoriel du banc de test thermodynamique en cours de finalisation}
  \caption{Schéma de principe du dispositif de mesure (support réservé)}
  \label{fig:emplacement_reserve}
\end{figure}

L'emplacement réservé à la~\cref{fig:emplacement_reserve} apparaît automatiquement avec un code couleur doux et s'enregistre dans le tableau récapitulatif des tâches.

\subsection{Bascule instantanée vers le mode publication finale}
\label{subsec:bascule_publication_finale}

Le cauchemar traditionnel de l'auteur sous \LaTeX{} est de devoir supprimer manuellement des dizaines de notes de relecture avant d'imprimer ou de soumettre le PDF définitif au jury.

\begin{boiteSynthese}{Passage en version finale en une seule manipulation}
  Dans \texttt{ClasseGuillaume}, le mode révision est entièrement piloté par une seule option de classe dans \texttt{main.tex} :
  \begin{itemize}
    \item \textbf{Phase de rédaction / relecture :} \texttt{\textbackslash documentclass[..., todoNote]\{ClasseGuillaume\}}
    \item \textbf{Version finale pour impression :} Il suffit de \textbf{retirer le mot-clé \texttt{todoNote}} de l'en-tête \texttt{\textbackslash documentclass}.
  \end{itemize}
  Instantanément, toutes les notes dans la marge s'évaporent, les surlignages colorés redeviennent du texte ordinaire rigoureusement immaculé, et le tableau des révisions disparaît du début du document. Aucune modification du texte source n'est requise !
\end{boiteSynthese}



\chapter{Référencement Croisé Moderne \& Bibliographie Académique}
\label{chap:references_biblio}
%------------------------------------------------------------------------------------
%  SECTION 8.1 : RÉFÉRENCES SÉMANTIQUES AVEC CLEVEREF
%------------------------------------------------------------------------------------
\section{Références Sémantiques Intelligentes avec \textsf{cleveref}}
\label{sec:references_cleveref}

Le renvoi vers des éléments numérotés (équations, figures, tableaux, sections, annexes) est au cœur de la rigueur scientifique.
En \LaTeX{} traditionnel, l'usage de la macro \texttt{\textbackslash ref\{label\}} est une source inépuisable de coquilles :
l'auteur doit taper manuellement le mot descriptif (« à la figure \texttt{\textbackslash ref\{...\}} », « dans le tableau \texttt{\textbackslash ref\{...\}} »).
En cas de remaniement, si un tableau devient une figure ou si le document est traduit, toutes ces mentions manuelles doivent être corrigées une à une.

Le modèle \texttt{ClasseGuillaume} intègre le paquet de référence \textsf{cleveref}, configuré pour la langue française avec les options \texttt{[french, capitalize, noabbrev]}.

\subsection{Les atouts majeurs de \texttt{\textbackslash cref} et \texttt{\textbackslash Cref}}
\label{subsec:atouts_cleveref}

\begin{itemize}
  \item \textbf{Détection sémantique automatique du type d'objet :}
  La commande \texttt{\textbackslash cref\{label\}} identifie automatiquement la nature de la cible et lui associe le mot idoine avec ses accords grammaticaux en français :
  \begin{itemize}
    \item Une équation devient automatiquement « \cref{eq:bernoulli_moderne} ».
    \item Un tableau devient « \cref{tab:synthese_mesures} ».
    \item Une figure devient « \cref{fig:schema_bloc_tikz} ».
    \item Une illustration personnalisée devient « \cref{ill:microscopie_nanotubes} ».
    \item Une section devient « \cref{sec:architecture_projet} ».
  \end{itemize}

  \item \textbf{Majuscule automatique en début de phrase (\texttt{\textbackslash Cref}) :}
  En tête de phrase, la macro majuscule \texttt{\textbackslash Cref\{label\}} applique la lettre capitale requise :
  \Cref{tab:comparatif_classe} récapitule les divergences fondamentales entre les approches.

  \item \textbf{Regroupement automatique des références multiples :}
  Lorsque plusieurs objets similaires sont appelés simultanément, \textsf{cleveref} les regroupe élégamment au pluriel :
  Les~\cref{eq:maxwell_faraday,eq:maxwell_ampere} décrivent les lois de l'électromagnétisme classique.
  Les~\cref{tab:synthese_mesures,tab:comparatif_classe} confrontent les métriques d'analyse.

  \item \textbf{Renvois de pagination sémantiques (\texttt{\textbackslash cpageref}) :}
  Pour situer physiquement un résultat dans le mémoire, la macro \texttt{\textbackslash cpageref\{label\}} génère le renvoi :
  le schéma fonctionnel est consultable à la~\cpageref{fig:schema_bloc_tikz}.
\end{itemize}

\begin{boiteSynthese}{Prise en charge du flottant personnalisé \texttt{illustration}}
  Pour que le nouveau flottant iconographique soit parfaitement géré, la classe déclare explicitement :
  \begin{verbatim}
  \crefname{illustration}{illustration}{illustrations}
  \Crefname{illustration}{Illustration}{Illustrations}
  \end{verbatim}
  Cette harmonisation garantit qu'aucune rustine manuelle n'est requise lors de la rédaction.
\end{boiteSynthese}


%------------------------------------------------------------------------------------
%  SECTION 8.2 : GESTION BIBLIOGRAPHIQUE AVEC BIBLATEX ET STYLE VANCOUVER
%------------------------------------------------------------------------------------
\section{Gestion Bibliographique Moderne et Style Vancouver}
\label{sec:bibliographie_vancouver}

La qualité des sources et leur référencement rigoureux constituent le socle de toute démarche scientifique.
L'ancien moteur \textsf{Bib\TeX{}} (1985), limité aux caractères ASCII 7 bits et dépendant de fichiers de style rigides (\texttt{.bst}), est aujourd'hui totalement dépassé.
\texttt{ClasseGuillaume} s'appuie sur le couple de référence contemporain : le paquet \textsf{biblatex} propulsé par le moteur de tri Unicode \textbf{Biber}.

\subsection{Indépendance absolue entre données brutes (\texttt{.bib}) et style de rendu}
\label{subsec:independance_bib_style}

L'un des atouts fondamentaux de l'écosystème \LaTeX{} réside dans la séparation étanche entre le fond et la forme :
\begin{itemize}
  \item Le fichier \texttt{bib/references.bib} constitue une \textbf{base de données brute et universelle}.
  Il est totalement \emph{agnostique} du style de publication et ne stocke que les métadonnées sémantiques (auteurs, titres, volumes, éditeurs, dates, identifiants DOI et URL).
  \item Le \textbf{style de restitution visuelle} (Vancouver, IEEE, APA, Chicago, Nature, etc.) est quant à lui gouverné de manière unique dans le document \LaTeX{} (option \texttt{style=vancouver} de \textsf{biblatex} dans la classe ou le préambule).
\end{itemize}

Ainsi, pour soumettre un même mémoire ou article à une revue exigeant le style auteur-date (ex. APA) plutôt qu'un style numérique séquentiel (Vancouver ou IEEE), il suffit de modifier un seul mot dans la configuration de \textsf{biblatex}, sans jamais devoir réécrire ou modifier la moindre ligne du fichier \texttt{references.bib}.

\subsection{Pourquoi le style Vancouver médical et technique ?}
\label{subsec:pourquoi_vancouver}

Dans les disciplines biomédicales, les sciences appliquées et l'ingénierie, le style de citation séquentiel numéroté dit \textbf{Style Vancouver} est la norme internationale préconisée par l'\emph{International Committee of Medical Journal Editors} (ICMJE).
Contrairement aux styles auteur-date (type APA ou Harvard) qui alourdissent les phrases par de longues parenthèses, le style Vancouver insère un numéro de référence sobre et compact entre crochets, sans interrompre la lecture.

\subsection{Diversité des modes de citation en contexte}
\label{subsec:modes_citations_contexte}

Selon la tournure syntaxique de la phrase, plusieurs macros de citation sont mises en œuvre :

\begin{itemize}
  \item \textbf{Citation narrative en position de sujet (\texttt{\textbackslash textcite}) :}
  Comme l'a démontré \textcite{knuth1984texbook} dans ses traités fondateurs, la précision typographique relève d'une exigence algorithmique.
  De même, \textcite{lamport1994latex} puis \textcite{mittelbach2004companion} ont démocratisé la séparation entre forme et contenu.
  
  \item \textbf{Citation factuelle globale (\texttt{\textbackslash cite} ou \texttt{\textbackslash parencite}) :}
  L'agencement des proportions de la page découle des règles de l'art typographique classique~\cite{bringhurst2004elements}.
  La suite KOMA-Script apporte une flexibilité inégalée pour les ouvrages complexes~\cite{kohm2024koma},
  tandis que \textsf{fontspec} régit le choix des polices OpenType sous Lua\LaTeX{}~\cite{robertson2020fontspec}.
  
  \item \textbf{Études médicales et essais cliniques :}
  Dans le domaine de la santé, le diagnostic des pathologies gynécologiques s'appuie sur des méta-analyses rigoureuses~\cite{kanti2024ultrasound} et l'évaluation prospective de biomarqueurs sanguins~\cite{burghaus2024}.
  Par ailleurs, l'impact des directives cliniques sur les cohortes observationnelles a été modélisé récemment~\cite{nieva2024impact}.

  \item \textbf{Thèses de doctorat et manuels de référence :}
  Les problématiques d'optimisation documentaire sous Lua\LaTeX{} sont approfondies dans la thèse de doctorat de \textcite{winant2026thesis}.
  Le manuel officiel de référence Lua\TeX{} détaille quant à lui l'ensemble des primitives de bas niveau~\cite{luatex2024manual}.
\end{itemize}

\subsection{Francisation soignée et liens de retour \texttt{backref}}
\label{subsec:francisation_backref}

La classe personnalise le moteur \textsf{biblatex} pour une intégration francophone sans faille :
\begin{itemize}
  \item Les dates de consultation des ressources numériques affichent automatiquement « [consulté le \dots] » à partir du champ \texttt{urldate}.
  \item Les adresses web affichent la mention « disponible sur : » suivie du lien cliquable césuré proprement par \textsf{xurl}.
  \item \textbf{La fonctionnalité majeure \texttt{backref = true} :}
  Dans la section bibliographique en fin de document, chaque référence comporte automatiquement \textbf{la liste cliquable des numéros de pages} où elle a été citée dans le corps de l'ouvrage !
  Cette fonctionnalité permet aux membres du jury de vérifier instantanément dans quel contexte précis chaque source a été mobilisée.
\end{itemize}

Les citations de texte s'appuient sur le paquet \textsf{csquotes} avec la commande \texttt{\textbackslash enquote\{...\}}, produisant automatiquement des guillemets français dits en chevrons («~\dots~») avec espaces insécables normées.



%------------------------------------------------------------------------------------
%  ANNEXES (Numérotation alphabétique avec préfixe "Annexe")
%------------------------------------------------------------------------------------
\appendix
%------------------------------------------------------------------------------------
%  ANNEXE A : AIDE-MÉMOIRE COMPLET DES COMMANDES DE CLASSEGUILLAUME.CLS
%------------------------------------------------------------------------------------
\chapter{Aide-Mémoire Exhaustif des Macros et Environnements}
\label{annexe:cheatsheet_macros}

Cette annexe récapitule sous forme de tableau de référence synthétique la totalité des commandes, environnements, styles et options introduits par la classe \texttt{ClasseGuillaume.cls}.

\begin{small}
\setlength{\tabcolsep}{3pt}
\begin{longtable}{p{3.9cm} p{4.6cm} p{6.8cm}}
  \caption{Aide-mémoire des fonctionnalités fournies par \texttt{ClasseGuillaume.cls}}
  \label{tab:cheatsheet_macros} \\
  \toprule
  \textbf{Commande / Option} & \textbf{Syntaxe et Arguments} & \textbf{Rôle et Effet Typographique} \\
  \midrule
  \endfirsthead

  \multicolumn{3}{c}{\tablename\ \thetable\ -- \textit{Suite de l'aide-mémoire}} \\
  \toprule
  \textbf{Commande / Option} & \textbf{Syntaxe et Arguments} & \textbf{Rôle et Effet Typographique} \\
  \midrule
  \endhead

  \midrule
  \multicolumn{3}{r}{\textit{Suite sur la page suivante \dots}} \\
  \bottomrule
  \endfoot

  \bottomrule
  \endlastfoot

  \multicolumn{3}{l}{\textbf{1. Options de Classe (\texttt{\textbackslash documentclass[...]\{ClasseGuillaume\}})}} \\
  \addlinespace
  \texttt{todoNote} & Mot-clé dans les options & Active l'affichage des notes marginales, surlignages et de la table des révisions. \\
  \texttt{font=calibri} & Option de police (par défaut) & Charge Calibri depuis le système Windows ou le dossier local de repli. \\
  \texttt{font=heros} & Option de police & Charge TeX Gyre Heros (équivalent libre d'Helvetica de \TeX{} Live). \\
  \texttt{font=libertinus} & Option de police & Charge Libertinus Serif (police académique avec empattements). \\
  \texttt{font=comicsans} & Option de police & Charge Comic Sans MS pour présentations pédagogiques. \\
  \texttt{font=standard} & Option de police & Charge la police standard Latin Modern. \\
  \addlinespace
  \midrule

  \multicolumn{3}{l}{\textbf{2. Commandes de Polices et Titrages Spécifiques}} \\
  \addlinespace
  \texttt{\textbackslash AsapFont} & \texttt{\{\textbackslash AsapFont Texte\}} & Active la police Asap Condensed (utilisée pour les têtes de chapitres). \\
  \texttt{\textbackslash CalibriFont} & \texttt{\{\textbackslash CalibriFont Texte\}} & Force localement la police sans-serif Calibri. \\
  \texttt{\textbackslash ComicSansMSFont} & \texttt{\{\textbackslash ComicSansMSFont Texte\}} & Force localement la police Comic Sans MS. \\
  \addlinespace
  \midrule

  \multicolumn{3}{l}{\textbf{3. Gestion Dynamique des Géométries de Page}} \\
  \addlinespace
  \texttt{\textbackslash geometrieWord} & Sans argument & Marges symétriques de \qty{2.5}{\centi\meter} et marge de notes de \qty{2.2}{\centi\meter}. \\
  \texttt{\textbackslash geometrieLarge} & Sans argument & Marges compactes de \qty{1.2}{\centi\meter} pour couverture, grands tableaux ou posters. \\
  \texttt{\textbackslash geometrieStandard} & Sans argument & Restaure le calcul de marges canonique KOMA-Script (\texttt{DIV=12}). \\
  \addlinespace
  \midrule

  \multicolumn{3}{l}{\textbf{4. Styles d'En-têtes et Pieds de Page (\textsf{scrlayer-scrpage})}} \\
  \addlinespace
  \texttt{scrheadings} & \texttt{\textbackslash pagestyle\{scrheadings\}} & Style par défaut avec filet haut (\qty{0.5}{pt}), section dynamique et pagination. \\
  \texttt{StyleFocus} & \texttt{\textbackslash pagestyle\{StyleFocus\}} & Style épuré sans filet, marque de section centrée et pagination en bas. \\
  \texttt{StyleTest} & \texttt{\textbackslash pagestyle\{StyleTest\}} & Style de débogage affichant \texttt{\textbackslash jobname.tex}, \texttt{LastPage} et badge orange. \\
  \addlinespace
  \midrule

  \multicolumn{3}{l}{\textbf{5. Calque de Filigrane Institutionnel / Logo (\textsf{scrlayer})}} \\
  \addlinespace
  \texttt{\textbackslash masquerLogo} & Sans argument & Masque le logo d'arrière-plan (sur page de garde et pages préliminaires). \\
  \texttt{\textbackslash afficherLogo} & Sans argument & Réactive l'affichage du logo dans le coin supérieur droit. \\
  \addlinespace
  \midrule

  \multicolumn{3}{l}{\textbf{6. Boîtes et Encadrés Graphiques (\textsf{tcolorbox})}} \\
  \addlinespace
  \texttt{cadreFigure} & \texttt{\textbackslash begin\{cadreFigure\}} & Boîte aux coins arrondis (\qty{10}{pt}) et bordure primaire pour sublimer schémas/TikZ. \\
  \texttt{boiteSynthese} & \texttt{\textbackslash begin\{boiteSynthese\}} & Boîte teintée (\texttt{primaryColor!6}) avec en-tête pour synthèses et théorèmes. \\
  \addlinespace
  \midrule

  \multicolumn{3}{l}{\textbf{7. Flottants Personnalisés et Listes de Documents}} \\
  \addlinespace
  \texttt{illustration} & \texttt{\textbackslash begin\{illustration\}} & Flottant dédié aux supports artistiques et photographiques distinct des figures. \\
  \texttt{\textbackslash listofillustrations} & Sans argument & Génère la table des illustrations en fin de document (via fichier \texttt{.loi}). \\
  \addlinespace
  \midrule

  \multicolumn{3}{l}{\textbf{8. Système Collaboratif Avancé (\textsf{todonotes} + \textsf{lua-ul})}} \\
  \addlinespace
  \texttt{\textbackslash todoH} & \texttt{\textbackslash todoH\{style\}\{coul\}\dots} & Surlignage résistant, troncature > 30 car., infobulle PDF et lien retour vers \texttt{\textbackslash listoftodos}. \\
  \texttt{\textbackslash listoftodos} & \texttt{\textbackslash listoftodos[Titre]} & Tableau récapitulatif de toutes les révisions (associé à \texttt{\textbackslash label\{table\_todos\}}). \\
  \texttt{\textbackslash missingfigure} & \texttt{\textbackslash missingfigure[opts]} & Réserve un emplacement graphique proportionné en cours de finalisation. \\
  \texttt{GuillaumeStyle} & Profil d'auteur & Stylise la note avec le nom d'auteur Guillaume et une ombre portée. \\
  \texttt{MarieStyle} & Profil d'auteur & Stylise la note pour les remarques de direction académique. \\
  \texttt{FloraStyle} & Profil d'auteur & Stylise la note pour les remarques méthodologiques. \\
\end{longtable}
\end{small}


%------------------------------------------------------------------------------------
%  ANNEXE B : GUIDE DES PALETTES DE COULEURS ET PERSONNALISATION ESTHÉTIQUE
%------------------------------------------------------------------------------------
\chapter{Nuanciers et Guide de Personnalisation Chromatique}
\label{annexe:guide_palettes}

L'identité visuelle d'un mémoire ou d'un rapport technique s'adapte au domaine d'étude ou à la charte institutionnelle de l'établissement d'attache.
Cette annexe fournit quatre nuanciers professionnels complets accompagnés de leur code source \LaTeX{} prêt à l'emploi.

\section{Code source des palettes thématiques recommandées}
\label{sec:code_palettes}

Pour appliquer l'une de ces ambiances à votre document, copiez simplement le bloc de redéfinition correspondant dans votre préambule (dans \texttt{main.tex}) juste après le chargement de la classe.

\subsection{Palette 1 : Océan Académique (Configuration par défaut)}
\begin{boiteSynthese}{Ambiance Océan Académique}
  \textcolor{NavyBlue}{\rule{12pt}{12pt}} \textbf{Primaire :} \texttt{NavyBlue} (Bleu marine profond)\\
  \textcolor{TealBlue}{\rule{12pt}{12pt}} \textbf{Secondaire :} \texttt{TealBlue} (Bleu sarcelle équilibré)\\
  \textcolor{Bittersweet}{\rule{12pt}{12pt}} \textbf{Accent :} \texttt{Bittersweet} (Terre cuite douce)\\
  \textcolor{red!50!black}{\rule{12pt}{12pt}} \textbf{Liens web :} \texttt{red!50!black} (Bordeaux sombre)
\end{boiteSynthese}

\subsection{Palette 2 : Prestige Humanités et Sciences Morales}
Idéale pour les mémoires d'histoire, de philosophie, de littérature ou de droit :
\begin{verbatim}
% Palette Prestige Humanités :
\colorlet{primaryColor}{Maroon}          % Rouge bordeaux universitaire
\colorlet{secondaryColor}{darkgray}        % Gris anthracite profond
\colorlet{accentColor}{BurntOrange}       % Ocre doré
\colorlet{mdtRed}{Maroon!80!black}        % Liens sobres
\end{verbatim}

\subsection{Palette 3 : Sciences de la Vie, Écologie et Agronomie}
Idéale pour la biologie végétale, la pharmacologie, la médecine environnementale et les sciences agronomiques :
\begin{verbatim}
% Palette Sciences du Vivant :
\colorlet{primaryColor}{ForestGreen}     % Vert sapin dense
\colorlet{secondaryColor}{CadetBlue}       % Vert d'eau satiné
\colorlet{accentColor}{RawSienna}        % Teinte terreuse
\colorlet{mdtRed}{ForestGreen!70!black}   % Liens vert sombre
\end{verbatim}

\subsection{Palette 4 : Technologies de l'Information et Aérospatiale}
Recommandée pour l'informatique théorique, le génie mécanique, les télécommunications et la physique des plasmas :
\begin{verbatim}
% Palette Technologies & Ingénierie :
\colorlet{primaryColor}{RoyalBlue}       % Bleu cobalt électrique
\colorlet{secondaryColor}{MidnightBlue}   % Bleu nuit rigoureux
\colorlet{accentColor}{Bittersweet}       % Orange vif d'attention
\colorlet{mdtRed}{RoyalBlue!80!black}     % Liens bleus soutenus
\end{verbatim}

\section{Définition de teintes hexadécimales sur-mesure (HTML / RGB)}
\label{sec:couleurs_hexadecimales}

Si vous devez vous conformer aux codes hexadécimaux précis de la charte graphique de votre université (par exemple la charte officielle UCLouvain ou le code couleur de votre entreprise), le paquet \textsf{xcolor} accepte la syntaxe HTML suivante :

\begin{verbatim}
% Définition précise par code hexadécimal :
\definecolor{BleuUCLouvain}{HTML}{002855}
\definecolor{OrUCLouvain}{HTML}{D4AF37}

% Affectation immédiate à ClasseGuillaume :
\colorlet{primaryColor}{BleuUCLouvain}
\colorlet{secondaryColor}{BleuUCLouvain!75!white}
\colorlet{accentColor}{OrUCLouvain}
\end{verbatim}
Cette méthode assure une harmonisation stricte avec les chartes graphiques les plus exigeantes sans aucune modification du fichier de classe.


%------------------------------------------------------------------------------------
%  ANNEXE C : GUIDE D'OUTILLAGE, COMPILATION ET BONNES PRATIQUES GIT
%------------------------------------------------------------------------------------
\chapter{Guide d'Outillage, Compilation et Bonnes Pratiques Git}
\label{annexe:guide_outillage}

La reproductibilité d'un travail scientifique impose une maîtrise complète des outils d'automatisation et de traçabilité.
Cette annexe détaille la chaîne de compilation et les conventions de versionnement recommandées pour le présent modèle.

\section{Le script d'automatisation PowerShell : \texttt{compile.ps1}}
\label{sec:script_compile}

Sous l'environnement Windows, les terminaux PowerShell utilisent historiquement une page de code OEM (CP 850 ou CP 1252) susceptible de dégrader l'affichage des caractères accentués dans les journaux de compilation.
Le script \texttt{scripts/compile.ps1} sécurise cette étape :

\begin{enumerate}
  \item Il sauvegarde la page de code d'entrée et de sortie de la console.
  \item Il bascule temporairement l'environnement en UTF-8 universel (Page de code 65001).
  \item Il déclenche l'utilitaire \texttt{latexmk} avec le moteur Lua\LaTeX{} et la redirection \texttt{-outdir=build}.
  \item Dès l'achèvement, il restaure la page de code d'origine de Windows sans altérer le reste de votre session.
\end{enumerate}

La commande d'exécution à la racine du projet est immédiate :
\begin{verbatim}
powershell -ExecutionPolicy Bypass -File scripts/compile.ps1
\end{verbatim}

\section{Nettoyage des fichiers intermédiaires : \texttt{clean.ps1}}
\label{sec:script_clean}

Durant les cycles de compilation, des dizaines de fichiers auxiliaires (\texttt{.aux}, \texttt{.bcf}, \texttt{.fls}, \texttt{.run.xml}, \texttt{.tdo}) s'accumulent dans le dossier \texttt{build/}.
Le script \texttt{scripts/clean.ps1} permet de purger ces résidus pour garantir une compilation vierge à froid sans supprimer vos images sources ni vos fichiers bibliographiques.

\section{Bonnes pratiques de versionnement sous Git}
\label{sec:bonnes_pratiques_git}

Un dépôt Git professionnel doit respecter des règles strictes de salubrité :
\begin{itemize}
  \item \textbf{Fichiers à versionner impérativement :}\\
  Tous les fichiers sources textuels (\texttt{*.tex}), la classe (\texttt{ClasseGuillaume.cls}),
  la base bibliographique (\texttt{references.bib}), les graphiques et images sources,
  les scripts PowerShell (\texttt{*.ps1}), et le fichier \texttt{.latexmkrc}.
  
  \item \textbf{Fichiers à ignorer formellement (\texttt{.gitignore}) :}
  La totalité du dossier \texttt{build/} et les extensions auxiliaires (\texttt{*.aux}, \texttt{*.log}, \texttt{*.toc}, \texttt{*.synctex.gz}, \texttt{*.pdf} intermédiaire).
  Seuls les fichiers sources doivent être conservés ; le document PDF final est un artefact dérivé qui se reconstruit de façon déterministe par simple compilation.

  \item \textbf{Diffs visuels avec \textsf{latexdiff} :}
  Grâce aux sauts de ligne sémantiques adoptés par ce modèle, vous pouvez utiliser l'utilitaire \texttt{latexdiff} pour générer automatiquement un document PDF de révision mettant en évidence en rouge les passages supprimés et en bleu les ajouts entre deux versions ou commits Git successifs.
\end{itemize}


%------------------------------------------------------------------------------------
%  ANNEXE D : SPÉCIMEN D'INCLUSION DE DOCUMENTS PDF EXTERNES AVEC PDFPAGES
%------------------------------------------------------------------------------------
\chapter{Spécimen d'Inclusion de Pièces Jointes avec \textsf{pdfpages}}
\label{annexe:specimen_externe}

Dans de nombreux livrables académiques, les annexes finales doivent rassembler des documents tiers au format PDF :
attestations de stage, fiches constructeur, procès-verbaux de comités d'éthique, ou formulaires administratifs officiels.

\section{Paramètres avancés de la commande \texttt{\textbackslash includepdf}}
\label{sec:parametres_pdfpages}

Le paquet \textsf{pdfpages} (intégré nativement dans \texttt{ClasseGuillaume}) offre une panoplie complète de réglages pour adapter l'insertion sans dénaturer la mise en page :

\begin{table}[htbp]
  \centering
  \small
  \caption{Options essentielles de la commande \texttt{\textbackslash includepdf}}
  \label{tab:options_pdfpages}
  \begin{tabularx}{\textwidth}{p{2.8cm} p{5.5cm} X}
    \toprule
    \textbf{Option} & \textbf{Valeurs types} & \textbf{Effet produit sur le document inséré} \\
    \midrule
    \texttt{pages} & \texttt{pages=\{1-3,5\}}, \texttt{pages=-\}} & Sélectionne les pages précises ou l'intégralité du fichier PDF source. \\
    \addlinespace
    \texttt{pagecommand} & \texttt{\{\textbackslash thispagestyle\{scrheadings\}\}} & Conserve les en-têtes et le numéro de page séquentiel du mémoire au lieu d'une page vierge. \\
    \addlinespace
    \texttt{scale} & \texttt{scale=0.90} & Réduit légèrement le document importé pour laisser respirer les marges et l'en-tête. \\
    \addlinespace
    \texttt{offset} & \texttt{offset=10mm -5mm} & Décale l'implantation horizontale et verticale pour corriger un centrage imparfait. \\
    \addlinespace
    \texttt{nup} & \texttt{nup=2x1}, \texttt{nup=2x2} & Imprime 2 ou 4 pages du document source sur une seule feuille A4 (format planche de contact). \\
    \addlinespace
    \texttt{landscape} & \texttt{landscape=true} & Force l'orientation paysage pour les schémas techniques ou plans d'ingénierie. \\
    \bottomrule
  \end{tabularx}
\end{table}

\section{Exemple type d'intégration de formulaire officiel}
\label{sec:exemple_type_pdfpages}

Le bloc de code suivant illustre l'insertion type d'un certificat d'agrément éthique multipage :

\begin{singlespace}
\begin{verbatim}
% Exemple d'inclusion complète avec conservation de la numérotation KOMA :
\includepdf[
  pages       = {1-2},
  scale       = 0.92,
  pagecommand = {\thispagestyle{scrheadings}},
  link        = true,
  linkname    = {annexe:certificat_ethique}
]{documents/certificat_comite_ethique.pdf}
\end{verbatim}
\end{singlespace}

Grâce à l'option \texttt{link = true}, le document importé s'inscrit directement dans l'arborescence des signets PDF (\emph{bookmarks}) et peut faire l'objet de renvois interactifs dans le corps de votre thèse.



%------------------------------------------------------------------------------------
%  POST-CONTENU (BACKMATTER : Listes de figures, tables et bibliographie)
%------------------------------------------------------------------------------------
\backmatter

\listoffigures
\listofillustrations
\listoftables

\printbibliography[heading=bibintoc, title={Références Bibliographiques}]

\end{document}
